\documentclass[11pt]{article}
\usepackage[a4paper,margin=22mm]{geometry}
\usepackage{fontspec}
\setmainfont{lmroman10-regular.otf}[BoldFont=lmroman10-bold.otf,ItalicFont=lmroman10-italic.otf,BoldItalicFont=lmroman10-bolditalic.otf]
\usepackage{amsmath,amssymb,amsthm,mathtools,graphicx,xcolor,hyperref,xurl}
\setlength{\emergencystretch}{4em}
\newtheorem{theo}{Theorem}[section]
\newtheorem{lemm}[theo]{Lemma}
\newtheorem{prop}[theo]{Proposition}
\newtheorem{coro}[theo]{Corollary}
\newtheorem{corol}[theo]{Corollary}
\newtheorem{conj}[theo]{Conjecture}
\newtheorem{Proposition}[theo]{Proposition}
\newtheorem{theoreme}[theo]{Theorem}
\theoremstyle{definition}
\newtheorem{defi}[theo]{Definition}
\newtheorem{exam}[theo]{Example}
\newtheorem{exem}[theo]{Example}
\newtheorem{rema}[theo]{Remark}
\newtheorem{nota}[theo]{Notation}
\newtheorem*{rema*}{Remark}
\newtheorem*{defi*}{Definition}
\newtheorem*{theo*}{Theorem}
\newtheorem*{lemm*}{Lemma}
\newtheorem*{prop*}{Proposition}
\newcommand{\enoncename}{}
\newtheorem{corpusenonce}[theo]{\enoncename}
\newenvironment{enonce}[1]{\renewcommand{\enoncename}{#1}\begin{corpusenonce}}{\end{corpusenonce}}
\newenvironment{enonce*}[1]{\par\medskip\noindent\textbf{#1.}\itshape}{\par\medskip}
\newenvironment{proofc}{\begin{proof}}{\end{proof}}
\newcommand{\Mc}{\mathcal{M}}
\newcommand{\pointir}{.\ ---\ }
\providecommand{\mathof}[1]{\mathrm{#1}}

\numberwithin{equation}{section}
\numberwithin{figure}{section}
\numberwithin{table}{section}
\def\endappendix{}
%~Mouliné par MaN_auto v.0.23.0 2020-06-18 09:32:49



\newfontfamily\CorpusBbm{Noto Sans Math}
\newcommand{\mathbbm}[1]{\mathord{\text{\CorpusBbm\char"1D7D9}}}
\usepackage{mathrsfs}

%%PACKAGES DESACTIVES
%%\usepackage{etoolbox}
%%\usepackage[retainorgcmds]{IEEEtrantools}
%%%%%%%%%%%%%%%%%%%%%%%%%%%%%


%\DeclareMathOperator{\sgn}{sgn}
%\DeclareMathOperator{\Id}{Id}
%\DeclareMathOperator{\dom}{dom}
%\DeclareMathOperator{\supp}{supp}
%\DeclareMathOperator{\argmin}{argmin}
%\DeclarePairedDelimiter\abs{\lvert}{\rvert}
%\DeclarePairedDelimiter\norm{\lVert}{\rVert}





%%\AtBeginDocument{%\newtheorem{myname}[theo]{My beautiful theorem} }
%\newcommand*\CDRcopyright{License CC BY SA 4.0 https://creativecommons.org/licenses/by-sa/4.0/}
%%


%
%%\providecommand{\email}[1]{\href{mailto:#1}{\nolinkurl{#1}}}
%\providecommand{\arxivref}[2]{\href{http://arxiv.org/abs/#1}{arXiv:#1}\ifblank{#2}{}{ [#2]} }

\newcommand\RR{\mathbb{R}}
\newcommand\Zz{\mathbf{Z}}
\newcommand\Pp{\mathcal{P}}
\newcommand\PP{\mathbb{P}}
\newcommand\EE{\mathbb{E}}
\newcommand\Indic[1]{\Ind_{\{#1\}}}
\newcommand\Ind{\mathbbm{1}}
\newcommand\from{\colon}
\newcommand\NN{\mathbb{N}}
\newcommand\FF{\mathcal{F}}
\newcommand\Uu{\mathcal{U}}
\newcommand\Gg{\mathcal{G}}
\newcommand\Ee{\mathcal{E}}
\newcommand\pp{\mathbf{p}}
\newcommand\e{\mathrm{e}}
\newcommand\Zzu{\Zz^{(u)}}

\newcommand{\crochet}[1]{{\langle #1 \rangle}}


\newcommand\dd{\mathof{d}}

\DeclarePairedDelimiterX\ip[2]{\langle}{\rangle}{#1,#2}

%\newcommand\iu{\mathof{i}}


%environnements à supprimer !! 
%%\newenvironment{eqnarr}{\begin{IEEEeqnarray}{rCl{\ignorespacesafterend}
%%\end{IEEEeqnarray}\ignorespacesafterend}

%%\newenvironment{eqnarr*}{\begin{IEEEeqnarray*}{rCl}}{\end{IEEEeqnarray*}\ignorespacesafterend}
%%\IEEEeqnarraydefcolsep{9}{4em}


%\newcommand{\eqnarrLHS}[1]{\IEEEeqnarraymulticol{3}{l}{#1} \\ \quad}
%\newcommand{\eqnarrLHSnn}[1]{\IEEEeqnarraymulticol{3}{l}{#1} \nonumber \\ \quad}
%\newcommand\dee{\partial}
%\newcommand\Dom{\mathcal{D}}

%\newcommand\nfrac[2]{#1/#2}
%\newcommand\sP{\mathrm{P}}
%\newcommand\sE{\mathrm{E}}
%\newcommand\sQ{\mathrm{Q}}
%\newcommand\downto{\downarrow}
%\newcommand\upto{\uparrow}

%\newcommand\Aa{\mathcal{A}}

%\newcommand\Hh{\mathcal{H}}






%%%%%%%%%%%%%%%%%%%%%%%%%%%%%%%%%%%%%%%%%%%%%%%%%%%%%%%%

\graphicspath{{./figures/}}

\newcommand*{\mk}{\mkern -1mu}
\newcommand*{\Mk}{\mkern -2mu}
\newcommand*{\mK}{\mkern 1mu}
\newcommand*{\MK}{\mkern 2mu}

\hypersetup{urlcolor=purple, linkcolor=blue, citecolor=red}

\newcommand*{\romanenumi}{\renewcommand*{\theenumi}{\roman{enumi}}}
\newcommand*{\Romanenumi}{\renewcommand*{\theenumi}{\Roman{enumi}}}
\newcommand*{\alphenumi}{\renewcommand*{\theenumi}{\alph{enumi}}}
\newcommand*{\Alphenumi}{\renewcommand*{\theenumi}{\Alph{enumi}}}
\let\oldtilde\tilde
\renewcommand*{\tilde}[1]{\mathchoice{\widetilde{#1}}{\widetilde{#1}}{\oldtilde{#1}}{\oldtilde{#1}}}
%%%%%%%%%%%%%%%%%%%%%%%%%%%%%%%%%%%%%%%%%%%%%%%%%%%%%%%%%

\begin{document}
\section*{\raggedright 2606: Strong Malthusian convergence in L1 for conservative growth-fragmentations with bounded fission and a strict boundary growth gap}
\noindent\textbf{Author:} Jean Bertoin; Alexander R. Watson\par\noindent\textbf{Source:} The strong Malthusian behavior of growth-fragmentation processes. Annales Henri Lebesgue 3 (2020), publisher version, DOI 10.5802/ahl.46. \url{https://ahl.centre-mersenne.org/item/AHL_2020__3__795_0/}
\par\noindent\textbf{License:} CC BY 4.0, \url{https://creativecommons.org/licenses/by/4.0/}. The original publisher PDF has the explicit license notice. See the saved source/license records.
\par\noindent\textbf{Editorial material note (not author text):} The selected problem is Theorem 1.1, strong Malthusian convergence in L1, under the exact conservative-partition, continuous bounded-fission, irreducible/aperiodic and strict boundary growth-gap hypotheses. The complete author article body retains the local Ulam--Harris construction, many-to-one formula, h-biased spine, second-moment, tightness, uniform-ergodicity and conditional-fluctuation arguments. Named return-transform background results remain explicit references; other statements are unselected context. All proof text and formulas below are editable author TeX. Mathematical wording, language and proof order are retained. Only the article class, title matter and bibliography presentation are adapted for independent compilation; no proof has been generated or repaired.
\par\noindent\textbf{Editorial difficulty:} H2: The proof constructs an h-biased spine and derives second-moment bounds for a potentially infinitely branching system, combines compact tightness with uniform ergodicity, and controls conditional fluctuations to obtain strong random-measure convergence. This branching-process and Feynman–Kac analysis exceeds ordinary qualifying examinations.
\par\noindent\textbf{Editorial retrieval and attribution:} Retrieved 2026-10-01. DOI 10.5802/ahl.46. Source attribution notice: ../licenses/2606-source-license.txt; complete legal text: ../licenses/CC-BY-4.0.txt. Original source excerpts and the existing source-material check are retained in ../sources/2606/.
\medskip\hrule\medskip
\noindent\textbf{Author statement, complete proof and local supporting text follow in source order.}
\section{Introduction}
This work is concerned with the large time asymptotic behavior of a class of branching Markov processes in continuous time, which we call \emph{growth-fragmentation processes}. These may be used to model the evolution of a population, for instance of bacteria, in which an individual reproduces by fission into two or more new individuals.

Each individual grows continuously, with the growth depending deterministically on the current mass of the individual, up to a random instant at which fission occurs. This individual, which may be thought of as a mother, is then replaced by a family of new individuals, referred to as her daughters. We assume that mass is preserved at fission, meaning that the mass of the mother immediately before the division is equal to the sum of the masses of her daughters immediately afterwards. The time at which the fission occurs and the masses of her daughters at fission are both random, and depend on the mass of the mother individual. After a fission event, the daughters are in turn viewed as mothers of future generations, and evolve according to the same dynamics, independently of the other individuals.

Mathematically, we represent this as a process in continuous time, $\Zz=(\Zz_t, t\geq 0)$, with values in the space of point measures on $(0,\infty)$. Each individual is represented as an atom in $\Zz_t$, whose location is the individual's mass. That is, if at time $t$ there are $n\in \NN \cup\{\infty\}$ individuals present, with masses $z_1,z_2,\dotsc$, then $\Zz_t = \sum_{i=1}^n \delta_{z_i}$, with $\delta_z$ the Dirac delta at $z \in (0,\infty)$.

Growth-fragmentation processes are members of the family of structured population models, which were first studied using analytic methods in the framework of linear integro-differential equations. To demonstrate this connection, consider the \emph{intensity measure} $\mu_t$ of $\Zz_t$, defined by $\crochet{\mu_t,f} = \EE[\crochet{ \Zz_t,f}]$ for all $f\in{\mathcal C}_c$. That is, $f$ is a continuous function on $(0,\infty)$ with compact support, and the notation $\crochet{m,f} = \int f \, \dd m$ is used for the integral of a function $f$ against a Radon measure $m$ on $(0,\infty)$, whenever this makes sense. In words, $\mu_t(A)$ describes the concentration of individuals at time $t$ with masses in the set $A\subset (0,\infty)$, and, informally, the evolution of the branching Markov process $\Zz$ entails that the family $(\mu_t)_{t\geq 0}$ solves an evolution equation (see~\cite{EN00} for background) of the form
\begin{equation}\label{E:GFE}
\frac{\dd }{\dd t} \crochet{\mu_t,f}= \crochet{ \mu_t, {\mathcal A} f},
\end{equation}
where the infinitesimal generator
\[
{\mathcal A}f(x) = c(x) f'(x) + {B}(x) \int_{\Pp} \left(\sum_{i=1}^{\infty}f(xp_i)-f(x)\right)\kappa(x, \dd \pp)
\]
is naturally associated to the dynamics of $\Zz$, $f$ is a smooth function in the domain of ${\mathcal A}$ and $\Pp = \{ \mathbf{p} = (p_1,p_2,\dotsc) : p_1 \ge p_2 \ge \dotsb \ge 0, \sum_{i=1}^\infty p_i = 1\}$. The meaning of this operator will be described precisely later, when we derive it in equation~\eqref{E:GFG}. Briefly, an individual of mass $x$ grows at speed $c(x)$, experiences fission at rate $B(x)$ and, if fission occurs, then the relative masses of the daughters are drawn from the distribution $\kappa(x,\cdot)$. We shall refer to~\eqref{E:GFE} as the \emph{growth-fragmentation equation}. We assume that
\begin{gather}
c\colon (0,\infty)\to (0,\infty) \text{ is a continuous function,}\label{e:assump-c}
\\
B\colon (0,\infty) \to [0,\infty) \text{ is a continuous bounded function, and}\label{e:assump-B}
\\
\kappa \text{ is a continuous probability kernel from } (0,\infty) \text{ to } \Pp.\label{e:assump-kappa}
\end{gather}

A fundamental problem in this analytic setting is to determine explicit conditions on the parameters governing the evolution of the system that ensure the so-called (asynchronous) \emph{Malthusian behavior}: for all $f\in{\mathcal C}_c$,
\begin{equation}\label{E:Malthus}
\EE[\crochet{ \Zz_t,f}]=\crochet{\mu_t,f} \sim \e^{\lambda t}\crochet{\mu_0,h } \crochet{\nu,f}\quad \text{as }t \to \infty,
\end{equation}
where $\lambda \in \RR$, $h$ is positive function, and $\nu$ a Radon measure on $(0,\infty)$ with $\crochet{\nu,h }=1$. When~\eqref{E:Malthus} holds, we call $\lambda$ the \emph{Malthus exponent} and $\nu$ the \emph{asymptotic profile}. There exists a vast literature on this topic, and we content ourselves here to cite a few contributions~\cite{BerGab2, CalvoDoumicPerthame, DDGW, Esco} amongst the most recent ones, in which many further references can be found.

Spectral analysis of the infinitesimal generator ${\mathcal A}$ often plays a key role for establishing~\eqref{E:Malthus}. Indeed, if there exist $\lambda \in \RR$, a positive function $h$ and a Radon measure $\nu$ that solve the eigenproblem
\begin{equation}\label{E:eigenp}
{\mathcal A}h =\lambda h \ ,\ {\mathcal A}'\nu =\lambda \nu\ ,\
\crochet{\nu,h }=1,
\end{equation}
with $\mathcal{A}'$ the adjoint operator to $\mathcal{A}$, then~\eqref{E:Malthus} follows rather directly. In this direction, the Perron--Frobenius paradigm, and more specifically the Krein--Rutman theorem (which requires compactness of certain operators related to ${\mathcal A}$) yield a powerful framework for establishing the existence of solutions to the eigenproblem~\eqref{E:eigenp}. Then $\lambda$ arises as the leading eigenvalue of ${\mathcal A}$, i.e., the eigenvalue with the maximal real part, and $h$ and $\nu$ respectively as a corresponding positive eigenfunction and dual eigenmeasure. This method has been widely used in the literature; see, for instance, \cite{BCG-fine,DoumGab, MS16,Per07}. We also point out the recent works~\cite{BanCloGab, BCGM-qs}; the first develops a generalization of Doeblin's conditions to deal with a non-homogeneous setting, and the second approaches~\eqref{E:Malthus} in a manner akin to the analysis of quasi-stationary distributions.

A stochastic approach for establishing~\eqref{E:Malthus}, which is based on the Feynman--Kac formula and circumvents spectral theory, has been developed by the authors in~\cite{ BeFK,BW} and Cavalli in~\cite{Cavalli}. To carry out this programme, we introduce, under the assumption
\begin{equation}\label{e:condc}
\sup_{x>0} c(x)/x<\infty,
\end{equation}
the unique strong Markov process $X$ on $(0,\infty)$ with generator
\[
\mathcal{G} f (x) = \frac{1}{x} \mathcal{A}\bar{f}(x) - \frac{c(x)}{x} f(x),
\]
where $\bar{f}(x) = x f(x)$. Assume that
\begin{equation}\label{e:assump}
\text{the Markov process $X$, with generator $\mathcal{G}$, is irreducible and aperiodic,}
\end{equation}
and define the Feynman--Kac weight
\[
\mathcal{E}_t = \exp\left(\int_0^t \frac{c(X_s)}{X_s} \, \dd s\right),
\]
and the Laplace transform
\[
L_{x,y}(q)
= \EE_x[\e^{- qH(y)} \mathcal{E}_{H(y)} \Indic{H(y)<\infty}],
\]
where $H(y)=\inf\{t>0: X_t=y\}$ denotes the first hitting time of $y$ by $X$. A weaker version of~\cite[Theorem~1.2] {BeFK} (see also~\cite[Theorem~1.1]{BW}) can then be stated as follows.

\setcounter{theo}{-1}
\begin{theo}\label{t:0}
Assume~\eqref{e:assump-c}, \eqref{e:assump-B}, \eqref{e:assump-kappa}, \eqref{e:condc} and~\eqref{e:assump}. Fix $x_0 > 0$. Define
\[
\lambda = \inf\{q\in \RR: L_{x_0,x_0}(q) < 1\}.
\]
Then, $\lambda$ is real and its value is independent of $x_0$. If
\begin{equation}\label{E:main}
\limsup _{x\to 0+} \frac{c(x)}{x}< \lambda \quad\text{and}\quad \limsup _{x\to \infty} \frac{c(x)}{x}<\lambda,
\end{equation}
then the Malthusian behavior~\eqref{E:Malthus} holds (so $\lambda$ is the Malthus exponent) with
\[
h(x) = xL_{x,x_0}(\lambda) \quad \text{and} \quad
\nu(\dd y) = \frac{\dd y}{h(y)c(y) \lvert L'_{y,y}(\lambda)\rvert}.
\]
\end{theo}

Indeed, in~\cite{BeFK}, it was even shown that~\eqref{E:main} implies that~\eqref{E:Malthus} occurs at exponential rate. Theorem~\ref{t:0} will form the basis of our work, the purpose of which is to investigate the analog of~\eqref{E:Malthus} for the random variable $\crochet{ \Zz_t,f}$ itself, rather than merely its expectation. More precisely, assuming for simplicity that the growth-fragmentation process $\Zz$ starts from a single individual with mass $x>0$ and writing $\PP_x$ for the corresponding probability law, we prove the following result:

\begin{theo}\label{t:main}
Under the assumptions of Theorem~\ref{t:0}, the process $\Zz$ exhibits \emph{strong Malthusian behavior}: for all $x>0$ and for $f$ any continuous function satisfying $\lVert f/h \rVert_{\infty}<\infty$, one has
\goodbreak
\begin{equation}\label{e:MalthusF}
\lim_{t\to \infty} \e^{-\lambda t} \crochet{ \Zz_t,f}= \crochet{\nu,f} W_{\infty} \quad \text{ in }L^1(\PP_x),
\end{equation}
where
\[
W_\infty=\lim_{t\to \infty} \e^{-\lambda t} \crochet{ \Zz_t,h} \quad\text{and} \quad \EE_x[W_{\infty}]=h(x).
\]
\end{theo}

A few remarks may be in order to explain~\eqref{E:main} informally. On the one hand, $c(x)/x$ represents the instantaneous growth rate for a mass $x>0$, whereas the Malthus exponent $\lambda$ rather describes the long-time average rate of growth of the total mass of the particle system. So~\eqref{E:main} requests the instantaneous growth rate to be smaller than the long time average growth rate, for particles close to the boundary. This can be interpreted as a condition guaranteeing that the main contribution to the growth of the system stems from particles located in some compact set in $(0,\infty)$. Roughly speaking, one expects the strong Malthusian behavior to hold for branching particle systems with a compact state space, and indeed, compactness also plays a crucial in the Krein--Rutman theory.

The criterion~\eqref{E:main} involves the Malthus exponent $\lambda$, which is itself usually not explicitly known. It might therefore appear unsatisfactory. However, one can easily obtain lower-bounds for $\lambda$ solely in terms of the characteristics of the growth-fragmentation process, and these yield a fully explicit criterion. We give an example of such a result as a conclusion to this work.

Of course, even though the Malthusian behavior~\eqref{E:Malthus} suggests that its strong version~\eqref{e:MalthusF} might hold, this is by no means automatic. For instance, it should be plain that~\eqref{e:MalthusF} cannot hold when $\lambda$ is negative. Furthermore, although one might expect that $W_\infty > 0$, $\PP_x$-a.s. (since plainly, $\Zz$ never becomes extinct under our assumptions), this is not obvious, and is known to fail in some other branching models~\cite{DB-bpre}. It appears to be difficult to prove this in our case, since the space of individual sizes is non-compact.

The question of strong Malthusian behavior has been considered in the literature on branching processes for several different models, including general Crump--Mode--Jagers branching processes~\cite{Jagers,JN84,Nerman}, branching random walks~\cite{Biggins92}, branching diffusions~\cite{BBHHR, EnHaKy, GiHaHa, HaHa, HaHeKy}, branching Markov processes~\cite{AsHer, CRY, ChYu, Shio} and pure fragmentation processes~\cite{Beresty,Ber-frag,BeRou}. A notable recent development is the study of the neutron transport equation and associated stochastic processes~\cite{CHHK, HHK,HKV}, which uses a different probabilistic approach based on the notion of quasi-stationarity, as in~\cite{ChVil,CV-unbounded}.

Focusing in a little more detail on the literature on growth-fragmentation processes, Bertoin et~al.~\cite{BBCK} and Dadoun~\cite{Dadoun} deal with the self-similar case $c(x) = a x^{\alpha+1}$, $B(x) = bx^{\alpha}$, $\kappa(x,\dd \pp) = \kappa(\dd \pp)$, and look at the convergence (in $L^p(\PP_x)$ for some $p>1$) of a reweighted collection of measures $\sum_{u}(Z^u_t)^{\omega}\delta_{t^{\frac{1}{\alpha}}Z^u_t}\Indic{b_u\le t<d_u}$, where $\omega$ is a particular weighting. Shi~\cite{ShiQ-ou} deals with the case $c(x) = (a - \theta \log x)x$, $B(x) = b$ and $\kappa(x,\dd \pp) = \kappa(\dd \pp)$ and proves convergence (in $\PP_x$-probability) of $\frac{1}{\#\mathbf{Z}_t} \sum_{u} \delta_{Z^u_t} \Indic{b_u\le t < d_u}$. These references in fact allow infinite-rate fragmentation, and in the case of~\cite{ShiQ-ou}, Gaussian fluctuations in the individuals' sizes are permitted. On the other hand, the coefficients are somewhat more specific there than our \linebreak assumptions allow. Cloez~\cite{Cloez} proves strong Malthusian behavior in a more general trait model, essentially assuming existence of a positive eigenfunction, an ergodicity result and some growth bounds, and verifies these in the case $c\equiv 1$, $B$ bounded away from $0$ and $\infty$ and $\kappa$ an atom at the point $\pp = (1/2,1/2,0,\dotsc)$. In a similar vein, Bansaye et~al.~\cite{BDMT} assumes $B$ constant, which permits the identification of a Feynman--Kac representation, and then proves strong Malthusian behavior under some assumptions of ergodicity. Marguet~\cite{Mar-lln} takes a different approach, considering a time-dependent branching rate $B(t,x)$. Under certain assumptions, which include Foster--Lyapunov-type conditions, \cite{Mar-lln} uses a time-dependent Feynman--Kac approach to prove strong Malthusian behavior, and verifies these assumptions in the case $c(x) = ax$, $B(t,x) =t\phi(x)$ for $\phi$ bounded away from $0$ and $\infty$, and $\kappa$ being the distribution of $(\Theta,1-\Theta,0,\dotsc)$, $\Theta$ uniform on $[\epsilon,1-\epsilon]$. Of course, these are just a sample of works on this topic, and many more references can be found cited within them.


Here, we view $(\Zz_t, t\ge 0)$ as a general branching process in the sense of Jagers \cite{Jagers}. This means that, rather than tracking the mass of individuals at a given time, we instead track the birth time, birth mass and death (i.e., fission) time of every individual in each successive generation; of course, since individuals grow deterministically until they split and die, this induces no loss of information. This process can be characterised in terms of a \emph{reproduction kernel}; given the birth time and mass of an individual, this describes the distribution of the birth times and masses of its daughters. Assuming that this general branching process is Malthusian and supercritical (as defined in~\cite[Section~5]{Jagers} in terms of the reproduction kernel), and that a certain $x\log x$ integrability condition and some further technical assumptions are fulfilled, \cite[Theorem~7.3]{Jagers} essentially states that~\eqref{e:MalthusF} holds with $W_{\infty}$ the terminal value of the so-called \emph{intrinsic martingale}. However, the assumptions and the quantities appearing in~\cite[Theorem~7.3]{Jagers} are defined in terms of the reproduction kernel, sometimes in an implicit way. It appears to be rather difficult to understand the hypotheses and conclusions of~\cite{Jagers} in terms of the parameters of the growth-fragmentation process; for instance, it does not seem to be straightforward to connect the general branching process with the eigenproblem~\eqref{E:eigenp}.



Our approach combines classical elements with some more recent ingredients. Given the Malthusian behavior recalled in Theorem~\ref{t:0}, the main technical issue is to find explicit conditions, in terms of the characteristics of the growth-fragmentation, which ensure the uniform integrability of a remarkable martingale that is closely related to the intrinsic martingale. More precisely, as Theorem~\ref{t:0} may suggest, we first establish a so-called many-to-one (or Feynman--Kac) formula, which provides an expression for the intensity measure $\mu_t$ of the point process $\Zz_t$ in terms of a functional of the (piecewise deterministic) Markov process $X$. Making use of results in~\cite{BW}, this enables us to confirm that $\mu_t$ indeed solves the growth-fragmentation equation~\eqref{E:GFE}, and to construct a remarkable additive martingale associated with the growth-fragmentation process $\Zz$, namely
\[
W_t=\e^{-\lambda t} \crochet{\Zz_t, h},\quad t\geq 0,
\]
where the Malthus exponent $\lambda$ and the function $h$ are defined in terms of the Markov process $X$. In fact, $W$ is nothing but the version in natural times of the intrinsic martingale indexed by generations, as defined in~\cite[Section~5]{Jagers}. We shall then prove that the boundedness in $L^2(\PP_x)$, and hence the uniform integrability, of the martingale $W$ follows from~\eqref{E:main} by adapting the well-known spinal decomposition technique (described in~\cite{BigKyp} for branching random walks) to our framework.

The spine process, which is naturally associated to the intrinsic martingale, plays an important role in the proof of the strong Malthusian behavior~\eqref{e:MalthusF}. Specifically, it yields a key tightness property for the random point measures $\Zz_t$, which then enables us to focus on individuals with masses bounded away from $0$ and from $\infty$. This is crucial to extend the original method of Nerman~\cite{Nerman} to our setting.

The rest of this paper is organized as follows. In Section~\ref{section2}, we describe the precise construction of the growth-fragmentation process $\Zz$, which is needed in Section~\ref{section3} to establish a useful many-to-one formula for the intensity measure $\mu_t$ of $\Zz_t$. In particular, a comparison with results in~\cite{BW} makes the connection with the growth-fragmentation equation~\eqref{E:GFE} rigorous. The $L^2$-boundedness of the intrinsic martingale is established in Section~\ref{section4} under the assumption~\eqref{E:main}, and we then prove the strong Malthusian behavior~\eqref{e:MalthusF} in Section~\ref{section5}. Section~\ref{section6} is devoted to providing explicit conditions on the characteristics of the growth-fragmentation that ensure~\eqref{E:main}.

\section{Construction of the growth-fragmentation process}\label{section2}
To start with, we introduce the three characteristics $c, {B}$ and $\kappa$ which govern the dynamics of the growth-fragmentation process. First, as described in the introduction, let $c\from (0,\infty) \to (0,\infty)$ be a continuous function satisfying assumption~\eqref{e:condc}, namely
\begin{equation*}
\sup_{x>0} c(x)/x<\infty.
\end{equation*}
$c$ describes the speed of growth of individuals as a function of their masses. For every $x_0>0$, the initial value problem
\begin{equation}\label{e:ode}
\begin{cases}
\dot{x}(t) &= c(x(t)), \quad t \geq 0, \\
x(0) &= x_0,
\end{cases}
\end{equation}
has a unique solution that we interpret as the mass at time $t$ of an individual with initial mass $x_0$ when no fission occurred before time $t$.

Next, we consider a bounded, continuous function ${B}\colon (0,\infty)\to [0,\infty)$, which specifies the rate at which a particle breaks (or branches) as a function of its mass. That is, the probability that no fission event has occurred by time $t>0$ when the mass at the initial time is $x_0$, is given by
\[
\PP_{x_0}[\text{no fission before time }t]=\exp\left(-\int_0^t {B}(x(s))\dd s\right)=\exp\left(-\int_{x_0}^{x(t)} \frac{{B}(y)}{c(y)}\dd y\right).
\]



To complete the description and specify the statistics at fission events, we need to introduce some further notation. We call a non-increasing sequence $\pp=(p_1, p_2, \ldots)$ in the unit sphere of $\ell^1$, i.e.,


\[
p_1\geq p_2\geq \dotsb \geq 0 \text{ and } \sum_{i\geq 1}p_i= 1,
\]
a (proper) \emph{mass partition}. In our setting, we interpret a mass partition as the sequence (ranked in the non-increasing order) of the daughter-to-mother mass ratios at a fission event, agreeing that $p_i=0$ when the mother begets less than $i$ daughters. The space $\Pp$ of mass partitions is naturally endowed with the $\ell^1$-distance and we write ${\mathcal B}(\Pp)$ for its Borel $\sigma$-algebra. We consider a continuous probability kernel
\[
\kappa \from (0,\infty)\times {\mathcal B}(\Pp) \to [0,1],
\]
in the sense that $x \mapsto \kappa(x,\cdot)$ is continuous with respect to weak convergence of probability measures. We think of $\kappa (x,\dd \pp)$ as the distribution of the random mass partition resulting from a fission event that occurs when the mother has mass $x>0$. We always implicitly assume that $\kappa(x, \dd \pp)$ has no atom at the trivial mass partition $(1,0,0, \ldots)$, as the latter corresponds to a fictive fission.


We next provide some details on the construction of growth-fragmentation processes and make the framework rigorous. We denote by $\Uu = \bigcup_{n\ge 0} \NN^n$ the Ulam--Harris tree of finite sequences of positive integers, which will serve as labels for the individuals. As usual, we interpret the length $|u|=n$ of a sequence $u\in\NN^n$ as a generation, and for $i\in\NN$, write $ui$ for the sequence in $\NN^{n+1}$ obtained by aggregating $i$ to $u$ as its $(n+1)^{\text{th}}$ element, viewing then $ui$ as the $i^{\text{th}}$ daughter of $u$. The unique element of $\NN^0$, written $\varnothing$, will represent an initial individual.

We fix $x_0>0$ and aim at constructing the growth-fragmentation process $(\Zz_t, t\geq 0)$ started from a single atom at $x_0$, which we understand to represent a single progenitor individual, \emph{Eve}. We denote by $\PP_{x_0}$ the corresponding probability measure. First consider a random variable $\zeta $ in $(0, \infty]$ with cumulative distribution function
\[
\PP_{x_0}[\zeta \leq t]= 1-\exp\left(-\int_{x_0}^{x(t)} \frac{{B}(y)}{c(y)}\dd y\right), \quad t\geq 0,
\]
where $x(\cdot)$ denotes the solution to the flow velocity~\eqref{e:ode} started from $x_0$. We view $\zeta $ as the fission time of Eve, and thus the trajectory of Eve is
\[
Z_t^{\varnothing}=x(t) \text{ for }t<\zeta.
\]
We further set $b^{\varnothing}= 0$ and $d^{\varnothing}=\zeta $, so $[b^{\varnothing}, d^{\varnothing})$ is the time interval during which Eve is alive. We also view $d^{\varnothing}$ as the birth-time of the daughters of Eve and thus set $b^i=d^{\varnothing}$ for every $i\in\NN$.
\goodbreak

Next, conditionally on $d^{\varnothing}=s<\infty$, that is, equivalently, on $Z_{d^{\varnothing}-}^{\varnothing}=x$ with $x=x(s)$, we pick a random mass partition $\pp=(p_1, \ldots)$ according to the law $\kappa(x,\dd \pp)$. We view $xp_1, xp_2, \ldots$ as the masses at birth of the daughters of Eve and continue the construction iteratively in an obvious way. That is, conditionally on $xp_i=y>0$, the lifetime $\zeta^i$ of the $i^{\text{th}}$ daughter of Eve has the same distribution as $\zeta $ under $\PP_{y}$. Further set $d^i=b^i+\zeta^i$, and the trajectory of the $i^{\text {th}}$ daughter of Eve is thus
\[
Z_{t}^{i}=x(t-b^i) \text{ for }t\in[b^i,d^i),
\]
with $x(\cdot)$ now denoting the solution to~\eqref{e:ode} started from $y$. We stress that, thanks to~\eqref{e:condc}, the boundary point $0$ is a trap for the flow velocity, in the sense that the solution to~\eqref{e:ode} with initial value $x(0)=0$ is $x(t)=0$ for all $t$. Thus $0$ serves a cemetery state for particles, and individuals with zero mass can be simply discarded.

This enables us construct recursively a trajectory $(Z^u_t: t\in[b^u, d^u))$ for every $u\in \Uu$, and the state of the growth-fragmentation at time $t$ is then given by the point measure on $(0,\infty)$ with atoms at the locations of the individuals alive at time $t$, viz.
\[
\Zz_t = \sum_{u\in\Uu} \Ind_{\{t\in[b^u, d^u)\}} \delta_{Z^u_t}.
\]
We stress that the number of individuals may explode at a finite time even in situations when every mother always begets finitely many children (see, e.g.~\cite{Savits}), and then infinitely many fission events may occur on any non-degenerate time interval. On the other hand, it is readily seen from our key assumption~\eqref{e:condc} that the total mass process increases at most exponentially fast, specifically
\[
\crochet{\Zz_t, {\mathrm {Id}}}\leq x\e^{\gamma t}, \quad \PP_x{\text{-a.s.}}
\]
where $\gamma=\sup_{x>0} c(x)/x$. Thus the point process $\Zz_t$ is always locally finite; however the growth-fragmentation is not always a continuous time Markov chain.


\section{A many-to-one formula}\label{section3}

The first cornerstone of our analysis is a useful expression for the expectation of the integral of some function with respect to the random point measure $\Zz_t$ in terms of a certain Markov process $X$ on $(0,\infty)$. In the literature, such identities are commonly referred to as many-to-one formulas, they go back to~\cite{KaP, Pey} and are known to play a crucial role in the analysis of branching processes.

%%A: ?? label au coeur du paragraphe, mettre la formule après 'probability kernel en display!!
Recall that a size-biased pick from a mass partition $\pp=(p_1, \ldots)$ refers to a random element $p_K$, where the distribution of the random index $K$ is $\PP(K=i)=p_i$ for $i\in\NN$. Size-biased picking enables us to map the probability kernel $\kappa$ on $(0,\infty)\times \Pp$\label{s:kbar}
into a kernel $\bar k$ on $(0,\infty)\times (0,1)$ by setting for every $x>0$
\[
\int_{(0,1)} g(r)\bar k(x,\dd r)= {B}(x) \int_{\Pp} \sum_{i=1}^{\infty} p_ig(p_i) \kappa(x, \dd \pp)
\]
for a generic measurable function $g\from (0,1)\to \RR_+$. We stress that $\int_{(0,1)} \bar k(x,\dd r)= {B}(x)$ since $\kappa$ is a probability kernel on the space of proper mass partitions. We then introduce the operator
\[
{\mathcal G} f (x) = c(x) f'(x) +\int_{(0,1)} (f(rx) - f(x)) \bar k(x, \dd r),
\]
say defined for functions $f: (0,\infty)\to \RR$ which are bounded and possess a bounded and continuous derivative. It is easily seen that ${\mathcal G}$ is the infinitesimal generator of a unique Markov process, say $X=(X_t, t\geq 0)$. Recall that we have assumed condition~\eqref{e:condc} and that $B$ is bounded. By a slight abuse, we also use the notation $\PP_{x_0}$ for the probability measure under which this piecewise deterministic Markov process starts from $X_0=x_0$.

The evolution of $X$ can be described in words as follows. The process is driven by the flow velocity~\eqref{e:ode} until it makes a first downwards jump; more precisely, the total rate of jump at state $x$ is $\int_{(0,1)}\bar k(x,\dd r)={B}(x)$. Further, conditionally on the event that a jump occurs when the process is about to reach $x$, the position after the jump is distributed according to the image of the probability law ${B}(x)^{-1}\bar k(x, \dd r)$ by the dilation $r\mapsto rx$. An alternative formulation which makes the connection to the growth-fragmentation process more transparent, is that $X$ follows the path of Eve up to its fission, then picks a daughter at random according to a size-biased sampling and follows her path, and so on, and so forth.

We now state a useful representation of the intensity measure of $\Zz_t$ in terms of the Markov process $X$.


\begin{lemm}[Many-to-one formula -- Feynman--Kac representation] \label{L1}\ \\*[0.10em]
Assume~\eqref{e:assump-c}--\eqref{e:assump-kappa} and~\eqref{e:condc}. Define, for every $t\geq 0$,
\[
\Ee_t = \exp\biggl\{\int_0^t \frac{c(X_s)}{X_s} \, \dd s\biggr\}.
\]
For every measurable $f\from (0,\infty)\to \RR_+$ and every $x_0>0$, we have
\[
\EE_{x_0} \left[\langle \Zz_t, f\rangle \right] = x_0\EE_{x_0}\left[\frac{f(X_t)}{X_t}\Ee_t \right].
\]
\end{lemm}

Lemma~\ref{L1} is closely related to~\cite[Lemma~2.2]{BW}, which provides a representation of the solution to the growth-fragmentation equation~\eqref{E:GFE} by Feynman--Kac formula. Specifically, introduce the growth-fragmentation operator ${\mathcal A}$ given for every $f\in{\mathcal C}^1_c$ by
\begin{equation}\label{E:GFG}
\begin{aligned}
{\mathcal A}f(x) &=\quad c(x) f'(x) + \int_{(0,1)} r^{-1} f(rx) \bar k(x, \dd r) -{B}(x) f(x)  \\
&=\quad c(x) f'(x) + {B}(x) \int_{\Pp} \left(\sum_{i=1}^{\infty}f(xp_i)-f(x)\right)\kappa(x, \dd \pp),
\end{aligned}
\end{equation}
then comparing Lemma~\ref{L1} above and~\cite[Lemma~2.2]{BW} shows that the intensity measure $\mu_t$ of $\Zz_t$ solves~\eqref{E:GFE} with $\mu_0=\delta_{x_0}$. A fairly natural approach for establishing Lemma~\ref{L1} would be to argue first that the intensity measure of $\Zz_t$ solves the growth-fragmentation equation for ${\mathcal A}$ given by~\eqref{E:GFG} and then invoke in~\cite[Lemma~2.2]{BW}. This idea is easy to implement when the number of daughters after a fission event is bounded (for instance, when fissions are always binary); however, making this analytic approach fully rigorous in the general case would be rather tedious, as the total number of individuals may explode in finite time and thus fission events accumulate. We rather follow a classical probabilistic approach and refer to the treatise by Del Moral~\cite{Delmoral} and the lecture notes of Shi~\cite{Shi} for background.

\begin{proof}
We set $T_0=0$ and then write $T_1<T_2<\dotsb$ for the sequence of the jump times of the piecewise deterministic Markov process $X$. We claim that for every generation $n\geq 0$, there is the identity
\begin{equation}\label{e:m21gen}
\EE_{x_0} \left[\sum_{|u|=n} f(Z^u_t) \Ind_{\{b^u \leq t < d^u\}} \right] = x_0\EE_{x_0}\left[\Ind_{\{T_n \leq t < T_{n+1}\}}\frac{f(X_t)}{X_t}\Ee_t \right].
\end{equation}
The many-to-one formula of Lemma~\ref{L1} then follows by summing over all generations.

We shall now establish~\eqref{e:m21gen} by iteration. The identity
\begin{equation}\label{e:Eexp}
\exp\left(\int_0^t \frac{c(x(s))}{x(s)}\dd s\right) = \frac{x(t)}{x(0)}
\end{equation}
for the solution to the flow velocity~\eqref{e:ode} makes~\eqref{e:m21gen} obvious for the generation $n=0$.


Next, by considering the fission rates of Eve, we get that for every measurable function $g\from [0,\infty)\times [0,\infty)\to \RR_+$ with $g(t,0)=0$, we have
\begin{multline}\label{e:interm1}
\EE_{x_0}\left[\sum_{i=1}^{\infty} g(b^i, Z^i_{b^i})\right] \\
=\int_0^{\infty} \dd t {B}(x(t)) \exp\left(-\int_0^t {B}(x(s)) \dd s\right) \int_{\Pp} \kappa(x(t), \dd \pp)
\sum_{i=1}^{\infty}g(t,x(t)p_i).
\end{multline}
We then write
\[
\sum_{i=1}^{\infty}g(t,x(t)p_i)=x(t) \sum_{i=1}^{\infty}p_i\frac{g(t,x(t)p_i)}{x(t)p_i},
\]
so that by comparing with the jump rates of $X$, we see that the right-hand side of~\eqref{e:interm1} equals
\[
\EE_{x_0}\left[\frac{g(T_1, X_{T_1})}{X_{T_1}}X_{T_1-}\right] = x_0 \EE_{x_0}\left[\frac{g(T_1, X_{T_1})}{X_{T_1}}\Ee_{T_1}\right],
\]
where we have used $\Ee_{T_1} = \Ee_{T_1-} = \frac{X_{T_1-}}{x_0}$ by~\eqref{e:Eexp}, since $(X_t: t < T_1)$ is a solution of~\eqref{e:ode}. Putting the pieces together, we have shown that
\begin{equation}\label{e:interm}
\EE_{x_0}\left[\sum_{i=1}^{\infty} g(b^i, Z^i_{b^i})\right]= x_0 \EE_{x_0}\left[\frac{g(T_1, X_{T_1})}{X_{T_1}}\Ee_{T_1}\right].
\end{equation}

We then assume that~\eqref{e:m21gen} holds for a given $n\geq 0$. Applying the branching property at the fission event of Eve, we get
\[
\EE_{x_0} \left[\sum_{|u|=n+1} f(Z^u_t) \Ind_{\{b^u \leq t < d^u\}} \right]=
\EE_{x_0}\left[\sum_{i=1}^{\infty} g(b^i, Z^i_{b^i})\right],
\]
with
\begin{equation}\label{e:interg}
g(s,y)= \EE_{y} \left[\sum_{|u|=n} f(Z^u_t) \Ind_{\{b^u \leq t-s < d^u\}} \right] = y\EE_y\left[\Ind_{\{T_n \leq t-s < T_{n+1}\}}\frac{f(X_{t-s})}{X_{t-s}}\Ee_{t-s}
\right]
\end{equation}
for $s\leq t$ and $g(s,y)=0$ otherwise, using the induction hypothesis~\eqref{e:m21gen} in the last equality. We conclude as follows:
\begin{align*}
x_0 \EE_{x_0}\left[\Indic{T_n\le t < T_{n+1}} \frac{f(X_t)}{X_t} \Ee_t \right] &= x_0 \EE_{x_0}\left[\Ee_{T_1} \frac{g(T_1,X_{T_1})}{X_{T_1}} \right]
\\
&=
\EE_{x_0}\left[\sum_{i=1}^\infty g(b^i,Z^i_{b^i}) \right]
\\
&=
\EE_{x_0}\left[\sum_{\lvert u \rvert = n+1} f(Z^u_t) \Indic{b^u\le t<d^u}\right],
\end{align*}
using in the first equality the Markov property and the multiplicative property $\Ee_{t} = \Ee_{T_1-t}\Ee_{T_1}$, in the second~\eqref{e:interm} and in the final equality~\eqref{e:interg}. This shows that the many-to-one formula~\eqref{e:m21gen} holds for the generation $n+1$. By induction, \eqref{e:m21gen} holds for any $n$.
\end{proof}

The many-to-one formula of Lemma~\ref{L1} connects the intensity measure of the branching process $\Zz$ to the instrumental Markov process $X$. We will use this to identify the martingale $(W_t, t\geq 0)$, whose terminal value $W_{\infty}$ plays a key role in strong Malthusian behavior~\eqref{e:MalthusF}. We will use this martingale in the next section to define a tilted measure $\tilde \PP$, under which both the process $\Zz$ and the size of a selected cell $\tilde{X}$ are defined simultaneously.

In the final section of this work, we shall also need a version of Lemma~\ref{L1} extended to the situation where, roughly speaking, individuals are frozen at times which are observable from their individual trajectories. Specifically, we define a \emph{simple stopping line} to be a functional $T$ on the space of piecewise continuous trajectories $z=(z_t)_{t\geq 0}$ and with values in $[0,\infty]$, such that for every $t\geq 0$ and every trajectory $z$, if $T(z)\leq t$, then $T(z)=T(z')$ for any trajectory $z'$ that coincides with $z$ on the time-interval $[0,t]$, and then we simply write $z_T=z_{T(z)}$. Typically, $T(z)$ may be the instant of the $j^{\text{th}}$ jump of $z$, or the first entrance time $T(z)=\inf\{t>0: z_t\in A\}$ in some measurable set $A\subset (0,\infty)$. The notion of a simple stopping line is a particular case of the more general stopping line introduced by Chauvin~\cite{Cha-stop}. The restriction simplifies the proofs somewhat, and will be sufficient for our applications later.

We next introduce the notion of ancestral trajectories. Recall from the preceding section the construction of the trajectory $(Z^u_t: t\in[b^u, d^u))$ for an individual\linebreak$u=u_1\ldots u_n\in\Uu$. The sequence of prefixes $u^j=u_1 \ldots u_j$ for $j=1, \ldots, n$ forms the ancestral lineage of that individual. Note that, as customary for many branching models, the death-time of a mother always coincides with the birth-time of her children, so every individual $u$ alive at time $t>0$ (i.e. with $b^u\leq t < d^u$) has a unique ancestor alive at time $s\in[0,t)$, which is the unique prefix $u^j$ with $b^{u^j}\leq s < d^{u^j}$. We can thus define unambiguously the mass at time $s$ of the unique ancestor of $u$ which is alive at that time, viz. $Z^u_s=Z^{u^j}_s$. This way, we extend $Z^u$ to $[0,d^u)$, and get the ancestral trajectory of the individual $u$.


For the sake of simplicity, for any simple stopping line $T$ and any trajectory $z$, we write $z_T=z_{T(z)}$, and define the point process of individuals frozen at $T$ as
\[
\Zz_T = \sum_{u\in\Uu} \Ind_{\{T(Z^u)\in[b^u, d^u)\}} \delta_{Z^u_T}.
\]
Note that $\Zz_T = 0$ if $T(Z^u) = \infty$ for all $u$.

\begin{lemm}\label{L1'}
Assume~\eqref{e:assump-c}--\eqref{e:assump-kappa} and~\eqref{e:condc}. Let $T$ be a simple stopping line. For every measurable $f\from (0,\infty)\to \RR_+$ and every $x_0>0$, we have
\[
\EE_{x_0} \left[\langle \Zz_T, f\rangle \right] = x_0\EE_{x_0}\left[\frac{f(X_T)}{X_T}\Ee_{T(X)}, T(X)<\infty \right].
\]
\end{lemm}

\begin{proof}
The proof is similar to that of Lemma~\ref{L1}, and we use the same notation as there.
In particular, we write $x(\cdot)$ for the solution to the flow velocity~\eqref{e:ode} started from $x(0)=x_0$, and set $T(x(\cdot))=t_0\in[0,\infty]$. By the definition of a simple stopping line, we have obviously that under $\PP_{x_0}$, $T(Z^{\varnothing}) = t_0$ a.s. on the event $0\leq T(Z^{\varnothing}) < d^{\varnothing}$, and also $T(X)=t_0$ a.s. on the event $0\leq T(X) < T_1$. Using~\eqref{e:Eexp}, we then get
\[
\EE \left[f(Z^{\varnothing}_T) \Ind_{\{b^{\varnothing} \leq T(Z^{\varnothing}) < d^{\varnothing}\}} \right] =x_0\EE_{x_0}\left[\Ind_{\{0 \leq T < T_{1}\}}\frac{f(X_T)}{X_T}\Ee_T \right].
\]
Just as in the proof of Lemma~\ref{L1}, it follows readily by induction that for every generation $n\geq 0$, there is the identity
\[
\EE_{x_0} \left[\sum_{|u|=n} f(Z^u_T) \Ind_{\{b^u \leq T(Z^u) < d^u\}} \right] = x_0\EE_{x_0}\left[\Ind_{\{T_n \leq T < T_{n+1}\}}\frac{f(X_T)}{X_T}\Ee_T \right],
\]
and we conclude the proof of Lemma~\ref{L1'} by summing over generations.
\end{proof}


\section{Boundedness of the intrinsic martingale in \texorpdfstring{$L^2(\PP)$}{L\textasciicircum 2(P)} }\label{section4}
In order to apply results from~\cite{BW,BeFK}, we shall now make some further fairly mild assumptions that will be enforced throughout the rest of this work. Specifically, we suppose henceforth that assumption~\eqref{e:assump} holds, namely that
\[
\text{the Markov process $X$, with generator $\Gg$, is irreducible and aperiodic.}
\]
Although~\eqref{e:assump} is expressed in terms of the Markov process $X$ rather than the characteristics of the growth-fragmentation process, it is easy to give some fairly general and simple conditions in terms of $c, B$ and $\kappa$ that guarantee~\eqref{e:assump}; see notably of~\cite[Lemma~3.1]{BeFK} for a discussion of irreducibility. We further stress that aperiodicity should not be taken to granted if we do not assume the jump kernel $\bar k$ to be absolutely continuous.

\begin{rema}
We mention that a further assumption is made in~\cite{BW,BeFK}, namely that the kernel $\bar k(x,\dd y)$ is absolutely continuous with respect to the Lebesgue measure, and that the function $(0,\infty) \ni x \mapsto \bar{k}(x,\cdot) \in L^1(0,\infty)$ is continuous. However, this is only needed in~\cite{BW} to ensure some analytic properties (typically, the Feller property of the semigroup, or the connection with the eigenproblem~\eqref{E:eigenp}), but had no role in the probabilistic arguments developed there.
We can safely drop this assumption here, and apply results of~\cite{BeFK,BW} for which it was irrelevant.
\end{rema}

Following~\cite{BW}, we introduce the Laplace transform
\[
L_{x,y}(q) = \EE_x\bigl[\e^{-qH(y)} \Ee_{H(y)} \Indic{H(y) < \infty} \bigr],\qquad q\in\RR,
\]
where $H(y) = \inf\{ t > 0: X_t = y \}$. For any $x_0>0$, the map $L_{x_0,x_0}\from \RR \to (0,\infty]$ is a convex non-increasing function with, $\lim_{q\to \infty}L_{x_0,x_0}(q)=0$ since $x\mapsto c(x)/x$ is bounded by assumption~\eqref{e:condc}. We then \emph{define} the \emph{Malthus exponent} as
\[
\lambda \coloneqq \inf\{q\in\RR: L_{x_0,x_0}(q)<1\}.
\]
Recall that the value of $\lambda$ does not depend on the choice for $x_0$, and that although our definition of the Malthus exponent apparently differs from that in~\cite[Section~5]{Jagers},  \cite[Proposition~3.3]{BW} strongly suggests that the two actually should yield the same quantity.


With this in place, we define the functions $\ell, h \from (0,\infty) \to (0,\infty)$ by
\[
\ell(x) = L_{x, x_0}(\lambda) \text{ and } h(x) = x \ell(x),
\]
and may now state the main result of this section.

\begin{theo}\label{T1}
Under the same assumptions as Theorem~\ref{t:main}, which are the conditions~\eqref{e:assump-c}--\eqref{e:assump-kappa}, \eqref{e:condc}--\eqref{e:assump} already developed together with~\eqref{E:main}, namely
\begin{equation*}
\limsup _{x\to 0+} \frac{c(x)}{x}< \lambda \quad\text{and}\quad \limsup _{x\to \infty} \frac{c(x)}{x}<\lambda,
\end{equation*}
for every $x>0$, the process
\[
W_t = \e^{-\lambda t} \crochet{\Zz_t, h}, \qquad t\geq 0
\]
is a martingale bounded in $L^2(\PP_x)$.
\end{theo}

Before tackling the core of the proof of Theorem~\ref{T1}, let us first recall some features proved in~\cite{BW,BeFK} and their immediate consequences. From~\cite[Section~3.5]{BeFK}, it is known that~\eqref{E:main} ensures the existence of some $q<\lambda$ with $L_{x_0,x_0}(q)<\infty$. By continuity and non-increase of the function $L_{x_0,x_0}$ on its domain, this guarantees that
\begin{equation}\label{e:assump:malthus}
L_{x_0, x_0}(\lambda)=1.
\end{equation}


\label{s:mart}
\cite[Theorem~4.4]{BW} then shows that $\e^{-\lambda t} \ell(X_t) \Ee_t$ is a $\PP_{x}$-martingale, and we can combine the many-to-one formula of Lemma~\ref{L1} and the branching property of $\Zz$ to conclude that $W_t$ is indeed a $\PP_x$-martingale. We therefore call $h$ a \emph{$\lambda$-harmonic function}; in this vein, recall also from~\cite[Corollary~4.5 and Lemma~4.6]{BW} that $h$ is an eigenfunction for the eigenvalue $\lambda$ of (an extension of) the growth-fragmentation operator ${\mathcal A}$ which has been defined in~\eqref{E:GFG}. We call $W=(W_t: t\geq 0)$ the \emph{intrinsic martingale}, as it bears a close connection to the process with the same name that has been defined in~\cite[Section~5]{Jagers}.


\goodbreak
\begin{rema}
This remark clarifies the connection with~\cite{Jagers}, and is not required for the remainder of the proofs.
\begin{enumerate}
\item \label{list4.3.1}It is convenient to view the atomic measure $\e^{-\lambda t} \Zz_t$ as a weighted version of point measure $\Zz_t$, where the weight of any individual at time $t$ is $\e^{-\lambda t}$. In this setting, $W_t$ is given by the integral of the $\lambda$-harmonic function $h$ with respect to the weighted atomic measure $\e^{-\lambda t} \Zz_t$. Next, consider for each $k\in \NN$, the simple stopping line $T_k$ at which a trajectory makes its $k^{\text {th}}$ jump, and recall from the preceding section, that $\Zz_{T_k}$ then denotes the point measure obtained from $\Zz$ by freezing individuals at the instant when their ancestral trajectories jump for the $k^{\text{th}}$ time. In other words, $\Zz_{T_k}$ is the point process that describes the position at birth of the individuals of the $k^{\text{th}}$ generation. Just, as above, we further discount the weight assigned to each individual at rate $\lambda$, so that the weight of an individual of the $k^{\text{th}}$ generation which is born at time $b$ is $\e^{-\lambda b}$ (of course, individuals at the same generation are born at different times, and thus have different weights). The integral, say ${\mathcal W}_k$, of the $\lambda$-harmonic function $h$ with respect to this atomic measure, is precisely the intrinsic martingale as defined in~\cite{Jagers}. Using more general stopping line techniques, one can check that the boundedness in $L^2$ of $(W_t: t\geq 0)$ can be transferred to $({\mathcal W}_k: k\in\NN)$. Details are left to the interested reader.

\item \label{list4.3.2}In fact, \eqref{e:assump:malthus}, which is a weaker assumption than~\eqref{E:main}, not only ensures that the process in continuous time $W_t=\e^{-\lambda t} \crochet{\Zz_t, h}$ is a martingale, but also that the same holds for the process indexed by generations, $({\mathcal W}_k: k\in\NN)$. Indeed, from the very definition of the function $L_{x_0,x_0}$, \eqref{e:assump:malthus} states that the expected value under $\PP_{x_0}$ of the nonnegative martingale $\e^{-\lambda t} \ell(X_t) \Ee_t$, evaluated at the first return time $H(x_0)$, equals $1$, and therefore the stopped martingale
\[
\e^{-\lambda t\wedge H(x_0)} \ell(X_{t\wedge H(x_0)}) \Ee_{t\wedge H(x_0)}, \quad t\geq 0
\]
is uniformly integrable. Plainly, the first jump time of $X$, $T_1$ occurs before $H(x_0)$, and the optional sampling theorem yields
\[
\EE_{x_0}[\e^{-\lambda T_1} \Ee_{T_1} \ell(X_{T_1})]=1.
\]
One concludes from the many-to-one formula of Lemma~\ref{L1'} (or rather, an easy pathwise extension of it) that $\EE_{x}[{\mathcal W}_1]=h(x)$ for all $x>0$, and the martingale property of ${\mathcal W}$ can now be seen from the branching property.
\end{enumerate}
\end{rema}


The rest of this section is devoted to the proof of Theorem~\ref{T1}; in particular we assume henceforth that~\eqref{E:main} is fulfilled.



To start with, we recall from~\cite[Lemma~4.6]{BW} that the function $\ell$ is bounded and continuous, and as a consequence
\begin{equation}\label{e:bornh}
\sup_{y>0} h(y)/y = \sup_{y>0} \ell(y)=\| \ell \|_{\infty}<\infty.
\end{equation}
Moreover $\ell$ and $h$ are strictly positive, and thus bounded away from $0$ on compact subsets of $(0,\infty)$. We shall use often these facts in the sequel.



The heart of the matter is thus to establish boundedness of $(W_t)_{t\ge 0}$ in $L^2(\PP_x)$, for which we follow the classical path based on the probability tilting and spine decomposition; see e.g.~\cite{BigKyp} and references therein.

Fix $t>0$. We define a probability measure $\tilde \PP_x$ on an augmented probability space by further distinguishing an individual $U_t$, called the spine, in such a way that
\[
\tilde{\PP}_x[\Lambda \cap \{ U_t = u \}] = h(x)^{-1}\e^{-\lambda t} \EE_x[{h}(Z^u_t) \Ind_{\Lambda}\Ind_{\{b^u \leq t < d^u\}}]
\]
for $\Lambda \in \FF_t = \sigma(\Zz_t, s \le t)$ and $u\in \Uu$. For $s\le t$, we define $U_s$ to be the ancestor at time $s$ of the individual $U_t$.

We then have that $\tilde{\PP}_x(\Lambda) = \EE_x\left[\Ind_{\Lambda} \cdot \frac{W_t}{W_0} \right]$, and the martingale property of $W$ ensures the consistency of the definition of $\tilde \PP_x$ with respect to the choice of $t$.

An equivalent description is to say that, under the conditional measure $ \tilde{\PP}_x[\cdot \mid \FF_t]$, the spine is picked at random amongst the individuals alive at time $t$ according to an $h$-biased sampling, and the ancestor $U_s$ of $U_t$ at time $s\leq t$ serves as spine at time~$s$.

In order to describe the dynamics of the mass of the spine $\tilde X_t=Z^{U_t}_t$ as time passes, we introduce first for every $x>0$
\[
w(x)=\int_{\Pp}\sum_{i=1}^{\infty} h(xp_i) \kappa(x, \dd \pp)
\]
and set
\begin{equation}\label{e:ratespine}
\tilde {B}(x)= \frac{w(x)}{h(x)}{B}(x)\quad\text{and}\quad \tilde \kappa(x, \dd \pp)= w(x)^{-1} \sum_{i=1}^{\infty} h(xp_i) \kappa(x, \dd \pp).
\end{equation}
In short, one readily checks that just as $X$, $\tilde X$ increases steadily and has only negative jumps. Its growth is driven by the flow velocity~\eqref{e:ode}, and the total rate of negative jumps at location $x$ is $\tilde {B}(x)$, which is the total fission rate of the spine when its mass is $x$. Further, $\tilde \kappa(x, \dd \pp)$ gives the distribution of the random mass partition resulting from the fission of the spine, given that the mass of the latter immediately before that fission event is $x$. At the fission event of the spine, a daughter is selected at random by $h$-biased sampling and becomes the new spine. We now gather some facts about the spine which will be useful later on.


\begin{lemm}\label{L20}
Assume conditions~\eqref{e:assump-c}--\eqref{e:assump-kappa}, \eqref{e:condc}--\eqref{E:main}. Let $\tilde X_t=Z^{U_t}_t$\linebreak represent the mass of the spine at time $t\geq 0$.
\begin{enumerate}
\item \label{list4.4.1}The law of the process $\tilde{X}$ can be expressed in terms of that of $X$ as follows: for $\Lambda \in \FF_t$,
\[
\tilde{\PP}_x(\Lambda) = \EE_x[\Ind_{\Lambda} \mathcal{M}_t],
\]
where $\mathcal{M}_t \coloneqq e^{-\lambda t} \Ee_t \ell(X_t)/\ell(X_0)$ is a $\PP_x$-martingale.

\item \label{list4.4.2}The process $\tilde X$ is Markovian and exponentially point recurrent, in the sense that if we write $\tilde H(y) =\inf\{t>0: \tilde X_t=y\}$ for the first hitting time of $y>0$ by $\tilde X$, then there exists $\varepsilon>0$ such that $\tilde{\EE}_x\bigl[\exp\bigl(\varepsilon \tilde H(y)\bigr)\bigr] <\infty$.

\item \label{list4.4.3}The following many-to one formula holds: for every nonnegative measurable function $f$ on $(0,\infty)$, we have
\begin{equation}\label{e:mtoY}
\EE_x[\crochet{\Zz_t, f}]=\e^{\lambda t} h(x)\tilde \EE_x\left[f\left(\tilde X_t\right)/h\left(\tilde X_t\right)\right].
\end{equation}

\item \label{list4.4.4}Any function $g\from (0,\infty)\to \RR$ such that $g\ell$ is continuously differentiable belongs to the domain of its extended infinitesimal generator $\tilde {\mathcal G}$ and
\begin{equation}\label{E:genY}
\tilde {\mathcal G}g(x) = \frac{1}{\ell(x)} {\mathcal G}(g\ell)(x) +(c(x)/x- \lambda) g(x),
\end{equation}
in the sense that the process
\[
g\left(\tilde X_t\right)-\int_0^t \tilde {\mathcal G}g\left(\tilde X_s\right) \, \dd s
\]
is a $\tilde \PP_x$-local martingale for every $x>0$.
\end{enumerate}
\end{lemm}
\begin{proofc}
\begin{enumerate}
\item \label{listproof4.4.1}It follows immediately from the definition of the spine and the many-to-one formula of Lemma~\ref{L1} that for every $t\geq 0$, the law of $\tilde X_t$ under $\tilde \PP_x$ is absolutely continuous with respect to that of $X_t$ under $\PP_x$, with density given by $\mathcal{M}_t = \e^{-\lambda t} \Ee_t \ell(X_t)/\ell(X_0)$, as claimed. We have already pointed out the martingale property of $\mathcal{M}$ above (see page~\pageref{s:mart}.)

\item \label{listproof4.4.2}In~\cite[Section~5]{BW}, we used $\mathcal{M}$ to construct a point-recurrent Markov $Y=(Y_t, t\geq 0)$ by probability tilting. Hence $Y$ has the same one-dimensional marginals as $\tilde X$, and since the two processes are Markovian, they have the same law. The claim that $\tilde X$ (that is equivalently, $Y$) is exponentially point recurrent then follows from~\cite[Proof of Theorem~2]{BeFK} and~\cite[Lemma~5.2$\MK$(iii)]{BW}.

\item \label{listproof4.4.3}The many-to-one formula~\eqref{e:mtoY} merely rephrases the very definition of the spine.

\item \label{listproof4.4.4}Finally, claim about the infinitesimal generator follows from \cite[Lem\-ma~5.1]{BW}.
\end{enumerate}
\end{proofc}

\begin{rema}
The description of the dynamics governing the evolution of the spine entails that its infinitesimal generator can also be expressed by
\begin{equation}\label{E:gentilde}
\tilde {\mathcal G}f(x) = c(x) f'(x) + \frac{B(x)}{h(x)} \int_{\Pp} \left(\sum_{i=1}^{\infty}h(xp_i)\bigl[f(xp_i)-f(x)\bigr]\right) \kappa(x, \dd \pp),
\end{equation}
say for any $f\in{\mathcal C}^1_c$. The agreement between~\eqref{E:genY} and~\eqref{E:gentilde} can be seen from the identity $\mathcal{G}\ell(x) = (\lambda - c(x)/x)\ell(x)$, which is proved in~\cite[Corollary~4.5$\MK$(i)]{BW}.
\end{rema}

We readily deduce from Lemma~\ref{L20} that the intensity measure of the growth-fragmentation satisfies the Malthusian behavior~\eqref{E:Malthus} uniformly on compact sets.

\begin{coro}\label{C1}
Assume conditions~\eqref{e:assump-c}--\eqref{e:assump-kappa}, \eqref{e:condc}--\eqref{E:main}. For every compact set $K\subset (0,\infty)$ and every continuous function $f$ with $\|f/h\|_{\infty}<\infty$, we have
\[
\lim_{t\to \infty} \e^{-\lambda t}\EE_x[\crochet{ \Zz_t,f}]= h(x) \crochet{\nu,f} \quad \text{uniformly for }x\in K,
\]
where the asymptotic profile is given by $\nu=h^{-1}\pi$, with $\pi$ the unique stationary law of the spine process $\tilde X$.
\end{coro}

\begin{proof}
Suppose $K\subset [b,b']$ for some $0<b<b'<\infty$, and fix $ \varepsilon \in (0,1)$.
For every $0<x<y$, let $s(x,y)$ denote the instant when the flow velocity~\eqref{e:ode} started from $x$ reaches $y$.

Since the total jump rate $\tilde B$ of $\tilde X$ remains bounded on $K$ and the speed of growth $c$ is bounded away from $0$ on $K$, we can find a finite sequence $b=x_1< x_2 < \ldots < x_j=b'$ such that for every $i=1, \ldots, j-1$ and every $x\in(x_i, x_{i+1})$:
\[
s(x,x_{i+1})<1 \quad \text{and}\quad s(x_i,x) < 1,
\]
as well as
\[
\tilde \PP_{x}\left(\tilde H(x_{i+1})=s(x,x_{i+1})\right) >1-\varepsilon \quad \text{and}\quad \tilde \PP_{x_i}\left(\tilde H(x)=s(x_i,x)\right) >1-\varepsilon.
\]
An immediate application of the simple Markov property now shows that for every $i=1, \ldots, j-1$, every $x\in(x_i, x_{i+1})$, every $t\geq 1$, and every nonnegative measurable function $g$, the following bounds hold
\[
(1-\varepsilon) \tilde \EE_{x_{i+1}}\left(g\left(\tilde X_{t-s(x,x_{i+1})}\right)\right)
\le
\tilde \EE_{x}\left(g\left(\tilde X_{t}\right)\right)
\le (1-\varepsilon)^{-1} \tilde \EE_{x_i}\left(g\left(\tilde X_{t+s(x_i,x)}\right)\right)
\]

On the other hand, we know that $\tilde X$ is irreducible, aperiodic and ergodic, with stationary law $\pi$ (recall Lemma~\ref{L20}$\MK$\eqref{listproof4.4.1}). Since $\varepsilon$ can be chosen arbitrarily small, it follows from above that $\tilde X$ is uniformly ergodic on $K$, in the sense that for every continuous and bounded function $g$,
\[
\lim_{t\to \infty} \tilde \EE_{x}\left(g\left(\tilde X_{t}\right)\right)= \crochet{\pi, g} \quad \text{uniformly for }x\in K.
\]
We conclude the proof with an appeal to the many-to-one formula of Lemma~\ref{L20}$\MK$\eqref{listproof4.4.2}, taking $g=f/h$.
\end{proof}


The next Lemma~\ref{L2} is a cornerstone of the proof of Theorem~\ref{T1}.

\begin{lemm}\label{L2}
Assume conditions~\eqref{e:assump-c}--\eqref{e:assump-kappa}, \eqref{e:condc}--\eqref{E:main}. We have:
\begin{enumerate}
\item \label{list4.7.1}There exists $a<\infty$ such that, for all $x>0$ and $t\geq 0$,
\[
\tilde \EE_x\left[1/\ell\left(\tilde X_t\right)\right]\leq at + 1/\ell(x).
\]
\item \label{list4.7.2}There exists some $\lambda'<\lambda$ such that, for all $x>0$,
\[
\lim_{t\to \infty} \e^{-\lambda' t} \tilde X_t=0\quad \text{ in }
L^{\infty}(\tilde \PP_x).
\]
\end{enumerate}
\end{lemm}

\begin{proofc}
\begin{enumerate}
\item \label{listproof4.7.1}We apply Lemma~\ref{L20}$\MK$\eqref{listproof4.4.4} to $g=1/\ell$, with
\[
\tilde {\mathcal G}\left(\frac{1}{\ell}\right) (x)= (c(x)/x-\lambda)/\ell(x).
\]
Our assumption~\eqref{E:main} ensures that the right-hand side above is negative for all $x$ aside from some compact subset of $(0,\infty)$. By taking the constant $a=\sup_{x>0} \tilde {\mathcal G}\left(1/{\ell}\right) <\infty$, we deduce from Lemma~\ref{L20}$\MK$\eqref{listproof4.4.3} by optional sampling that $\tilde \EE_x[1/\ell(\tilde X_t)] -at \leq 1/\ell(x)$, which entails our claim.

\item \label{listproof4.7.2}Recall from the description of the dynamics of the spine before the statement that $\tilde X$ increases continuously with velocity $c$ and has only negative jumps. As a consequence, $\tilde X$ is bounded from above by the solution to the flow velocity~\eqref{e:ode}. One readily deduces that $\lim_{t\to \infty} \e^{-\lambda' t}x(t)=0$ for every $\lambda' >\limsup_{x\to \infty} c(x)/x$, and since $\limsup_{x\to \infty} c(x)/x<\lambda$ according to our standing assumption~\eqref{E:main}, this establishes our claim.\qedhere
\end{enumerate}
\end{proofc}

\goodbreak

We now have all the ingredients needed to prove Theorem~\ref{T1}.

\begin{proof}[Proof of Theorem~\ref{T1}]
Since the projection on ${\mathcal F}_t$ of $\tilde \PP_x$ is absolutely continuous with density $W_t/W_0$, the process $W$ is bounded in $L^2(\PP_x)$ if and only if $
\sup_{t\geq 0} \tilde \EE_x [W_t] <\infty$. We already know that$\sup_{t\geq 0} \tilde \EE_x [\e^{-\lambda t} \tilde X_t] <\infty$, by~Lemma~\ref{L2}$\MK$\eqref{listproof4.7.2}, and we are thus left with checking that
\begin{equation}\label{e:goal}
\sup_{t\geq 0} \tilde \EE_x [W'_t] <\infty,
\end{equation}
where
\[
W'_t=W_t-\e^{-\lambda t} h\left(\tilde X_t\right).
\]

In this direction, it is well-known and easily checked that the law of $\Zz$ under the new probability measure $\tilde{\PP}_x$ can be constructed by the following procedure, known as the spine decomposition. After each fission event of the spine, all the daughters except the new spine start independent growth-fragmentation processes following the genuine dynamics of $\Zz$ under $\PP$. This spine decomposition enables us to estimate the conditional expectation of $W'_t$ under $\tilde \PP_x$, given the spine and its sibling. At each time, say $s>0$, at which a fission occurs for the spine, we write $\pp(s)=(p_1(s), \ldots)$ for the resulting mass partition, and $I(s)$ for the index of the daughter spine. Combining this with the fact that $(W_s: s\geq 0)$ is a $\PP_y$-martingale for all $y>0$ entails the identity
\[
\tilde \EE_x\left[W'_t\mid \left(\tilde X_s, \pp(s), I(s)\right)_{s\geq 0}\right] =
\sum_{s\in \tilde F, s\leq t} \e^{-\lambda s} \sum_{i\neq I(s)} h\left(\tilde X_{s-} p_i(s)\right)\,,
\]
where $\tilde F$ denotes the set of fission times of the spine. We see from~\eqref{e:bornh} that $\sum h(xp_i) \leq x \| \ell\|_{\infty}$ for every $x>0$ and every mass-partition ${\mathbf p}=(p_1, \ldots)$, so the right-hand side is bounded from above by
\[
\| \ell \|_{\infty} \sum_{s\in \tilde F, s\leq t} \e^{-\lambda s} \tilde X_{s-},
\]
and to prove~\eqref{e:goal}, we now only need to check that
\[
\tilde \EE_x \left[\sum_{s\in \tilde F} \e^{-\lambda s} \tilde X_{s-}\right] <\infty.
\]

Recall from~\eqref{e:ratespine} that $\tilde {B} = w{B}/h$ describes the fission rate of the spine, and observe from~\eqref{e:bornh} that $w(x)\leq \| \ell\|_{\infty} x$, so that
\[
\tilde {B}(x) \leq \| \ell\|_{\infty} \| {B} \|_{\infty}
\frac{1}{\ell(x)}
\qquad \text{for all }x>0.
\]
Using the fact that the predictable compensator of the fission times and relative offspring sizes $(s,\pp)$ is $\tilde{B}(\tilde{X}_{s-}) \kappa(\tilde{X}_{s-},\dd \pp)\,\dd s$, we obtain that
\[
\tilde \EE_x \left[\sum_{s\in \tilde F} \e^{-\lambda s} \tilde X_{s-}\right] \leq
\| \ell\|_{\infty} \| {B} \|_{\infty}\tilde \EE_x\left[\int_0^{\infty} \e^{-\lambda s} \frac{\tilde X_s}{\ell\left(\tilde X_s\right)}\dd s\right].
\]
We now see that the expectation in the right-hand side is indeed finite by writing first
\[
\int_0^{\infty} \e^{-\lambda s} \frac{\tilde X_s}{\ell\left(\tilde X_s\right)}\dd s
= \int_0^{\infty} \frac{1}{\ell\left(\tilde X_s\right)} \cdot \e^{-\lambda' s}\tilde X_s
\cdot \e^{-(\lambda-\lambda')s} \, \dd s
\]
and then applying Lemma~\ref{L2}.
\end{proof}



\section{Strong Malthusian behavior}\label{section5}

We assume again throughout this section that the assumption~\eqref{E:main} is fulfilled. We will prove Theorem~\ref{t:main}: the strong Malthusian behavior~\eqref{e:MalthusF} then holds.


The proof relies on a couple of technical lemmas. Recall from Section~\eqref{section3} the notation $Z^u\colon [0, d^u)\to(0,\infty)$ for the ancestral trajectory of the individual $u$.

The first Lemma~\ref{L:B} states a simple tightness result.

\begin{lemm}\label{L:B}
Assume conditions~\eqref{e:assump-c}--\eqref{e:assump-kappa}, \eqref{e:condc}--\eqref{E:main}. For every $x>0$ and $\varepsilon >0$, there exists a compact $K\subset (0,\infty)$ such that for all $t\geq 0$:
\[
\e^{-\lambda t} \EE_{x}\left[\sum_{u\in\Uu: b^u\leq t < d^u } h(Z^u_t) \Indic{Z^u_t\not\in K}
\right]< \varepsilon.
\]
\end{lemm}

\begin{proof}
From the very definition of the spine $\tilde X$, there is the identity
\[
\e^{-\lambda t} \EE_{x}\left[\sum_{u\in\Uu: b^u\leq t < d^u } h(Z^u_t) \Indic{Z^u_t\not\in K}
\right] = h(x) \tilde \PP_{x}\left[\tilde X_t\not\in K\right].
\]
Recall from Lemma~\ref{L20}$\MK$\eqref{listproof4.4.1} that $\tilde X$ is positive recurrent; as a consequence the family of its one-dimensional marginals under $\tilde \PP_{x}$ is tight, which entails our claim.
\end{proof}

The second Lemma~\ref{L:C} reinforces the boundedness in $L^2$ of the intrinsic martingale, cf.~Theorem~\ref{T1}.

\begin{lemm}\label{L:C}
Assume conditions~\eqref{e:assump-c}--\eqref{e:assump-kappa}, \eqref{e:condc}--\eqref{E:main}. For every compact subset $K\subset(0,\infty)$, we have
\[
\sup_{x\in K} \sup_{t\geq 0} \EE_x\left[W^2_t \right] < \infty.
\]
\end{lemm}

\begin{proof}
We may assume that $K=[b,b']$ is a bounded interval. For any $x\in(b,b']$, we write $s(x)$ for the time when the flow velocity~\eqref{e:ode} started from $b$ reaches $x$. We work under
$\PP_{b}$ and consider the event $\Lambda_x$ that the Eve individual hits $x$ before a fission event occurs.

We have on the one hand, that the law of $\Zz_{s(x)+t}$ conditionally on $ \Lambda_x$ is the same as that of $\Zz_t$ under $\PP_x$. In particular, the law of $W_t$ under $\PP_x$ is the same as that of $\e^{\lambda s(x)} W_{s(x)+t}$ under $\PP_b[\cdot \mid \Lambda_x]$, and thus
\[
\sup_{t\geq 0} \EE_x\left[W^2_t \right] \leq \frac{\e^{\lambda s(x)}}{\PP_b[\Lambda_x]} \EE_b\left[W^2_{\infty} \right].
\]

On the other hand, for every $x\in(b,b']$, we have $s(x)\leq s(b')<\infty$ and
\[
\PP_b[\Lambda_x]\geq \PP_b[\Lambda_{b'}]=\exp\left(-\int_{b}^{b'} \frac{{B}(y)}{c(y)}\dd y\right)>0,
\]
and our claim is proven.
\end{proof}

We have now all the ingredients to prove Theorem~\ref{t:main}.

\begin{proof}[Proof of Theorem~\ref{t:main}]
We suppose that $0\leq f \leq h$, which induces of course no loss of generality. Our aim is to check that $\e^{-\lambda(t+s)} \crochet{\Zz_{t+s},f} $ is arbitrarily close to $\crochet{\nu,f}W_t $ in $L^1(\PP_x)$ when $s$ and $t$ are sufficiently large. In this direction, recall that ${\mathcal F}_t$ denotes the natural filtration generated by $\Zz_t$ and use the branching property at time $t$ to express the former quantity as
\[
\e^{-\lambda (t+s)}\crochet{\Zz_{t+s},f} =
\sum_{u\in\Uu: b^u\leq t < d^u}
\e^{-\lambda t} h(Z^u_t) \cdot \frac{1}{h(Z^u_t)}\e^{-\lambda s}\crochet{\Zzu_{s},f},
\]
where conditionally on ${\mathcal F}_t$, the processes $\Zzu$ are independent versions of the growth-fragmentation $\Zz$ started from $Z^u_t$.

Fix $\varepsilon>0$. To start with, we choose a compact subset $K\subset (0,\infty)$ as in Lemma~\ref{L:B}, and restrict the sum above to individuals $u$ with $Z^u_t\not \in K$. Observe first that, since $\crochet{\Zzu_s, f}\leq \crochet{\Zzu_s, h}$ and $h$ is $\lambda$-harmonic, taking the conditional expectation given ${\mathcal F}_t$ yields
\[
\EE_x\!\left[\sum_{\substack{u\in\Uu, \\ b^u\leq t < d^u}}\e^{-\lambda t} h(Z^u_t) \Indic{Z^u_t\not \in K}\cdot \frac{1}{h(Z^u_t)}\e^{-\lambda s}\crochet{\Zzu_{s},f}\!\right]
\leq \EE_x\!\left[\sum_{\substack{u\in\Uu, \\ b^u\leq t < d^u}}\e^{-\lambda t} h(Z^u_t) \Indic{Z^u_t\not \in K}\!\right]\!.
\]


From the very choice of $K$, there is the bound
\begin{equation}\label{E:P1}
\EE_x\left[\sum_{u\in\Uu: b^u\leq t < d^u}\e^{-\lambda t} h(Z^u_t) \Indic{Z^u_t\not \in K}\cdot \frac{1}{h(Z^u_t)}\e^{-\lambda s}\crochet{\Zzu_{s},f}\right]\leq \varepsilon.
\end{equation}

Next, recall from Lemma~\ref{L:C} that
\[
C(K)\coloneqq
\sup_{y\in K} \sup_{s\geq 0} \EE_y\left[W_s^2 \right] < \infty,
\]
and consider
\[
A(u,t,s)=\frac{1}{h(Z^u_t)}\e^{-\lambda s}\crochet{\Zzu_{s},f}
\]
together with its conditional expectation given ${\mathcal F}_t$
\[
a(u,t,s)=\EE_x[A(u,t,s)\mid {\mathcal F}_t].
\]
Again, since $0\leq f \leq h$, for every $u$ with $Z^u_t\in K$, we have
\begin{equation}\label{e:inter-Aa}
\EE_x[(A(u,t,s)-a(u,t,s))^2\mid {\mathcal F}_t] \leq 4 C(K).
\end{equation}


Since conditionally on ${\mathcal F}_t$, the variables $A(u,t,s)-a(u,t,s)$ for $u\in{\mathcal U}$ are independent and centered, there is the identity
\begin{multline*}
\EE_x\left[\left| \sum_{u\in\Uu: b^u\leq t < d^u}\e^{-\lambda t} h(Z^u_t) \Indic{Z^u_t \in K}\cdot (A(u,t,s)-a(u,t,s))\right |^2
\right]\\
= \EE_x\left[\sum_{u\in\Uu: b^u\leq t < d^u}\e^{-2\lambda t} h^2(Z^u_t) \Indic{Z^u_t \in K}\cdot \EE_x\bigl[(A(u,t,s)-a(u,t,s))^2\bigm\vert {\mathcal F}_t\bigr]
\right],
\end{multline*}
and we deduce from~\eqref{e:inter-Aa} and the martingale property of $W$ that this quantity is bounded from above by
\[
4 C(K)\e^{-\lambda t} h(x) \max_{y\in K} h(y).
\]
This upper-bound tends to $0$ as $t\to \infty$, and it thus holds that
\[
\EE_x\left[\left| \sum_{u\in\Uu: b^u\leq t < d^u}\e^{-\lambda t} h(Z^u_t) \Indic{Z^u_t \in K}\cdot (A(u,t,s)-a(u,t,s))\right |
\right]< \varepsilon
\]
for all $t$ sufficiently large.

On the other hand, writing $y=Z^u_t$, we have from the branching property
\[
a(u,t,s)= \frac{1}{h(y)} \e^{-\lambda s} \EE_{y} \left[\crochet{\Zz_{s},f}\right],
\]
and Corollary~\ref{C1} entails that for all $s$ sufficiently large, $|a(u,t,s)-\crochet{\nu,f}|\leq \varepsilon$ for all individuals $u$ with $Z^u_t\in K$. Using the bound~\eqref{E:P1} with $h$ in place of $f$ for individuals $u$ with $Z^u_t\not \in K$ and putting the pieces together, we have shown that for all $s, t$ sufficiently large,
\[
\EE_x\left[\left| \e^{-\lambda (t+s)}\crochet{\Zz_{t+s},f} - \crochet{\nu,f} W_t
\right | \right] \leq (2+h(x))\varepsilon,
\]
which completes the proof of Theorem~\ref{t:main}.
\end{proof}


\section{Explicit conditions for the strong Malthusian behavior}\label{section6}

The key condition for strong Malthusian behavior, \eqref{E:main}, is given in terms of the Malthus exponent $\lambda$, which is not known explicitly in general. In this final section, we recall some known results that can be used to verify~\eqref{E:main}, and prove a sharper new one.

Recently, the following result was pointed out in~\cite{BeFK}; we provide a short outline of the proof.



\begin{prop}\label{P0}
Assume conditions~\eqref{e:assump-c}--\eqref{e:assump-kappa}, \eqref{e:condc}--\eqref{e:assump}, that
\[
d \coloneqq \inf_{x>0} \frac{c(x)}{x} =\lim_{x\to 0+} \frac{c(x)}{x} = \lim_{x\to \infty} \frac{c(x)}{x} \geq 0,
\]
and that $c$ is not the linear function $x \mapsto dx$. Assume moreover that there exist some $q_{\infty}>0$ and $x_{\infty}>0$ such that
\begin{equation}\label{e:lyap1a}
q_{\infty} c(x)/x + \int_{(0,1)} (r^{q_{\infty}}-1)\bar k(x,\dd r) \leq 0 \quad\text{for all }x\geq x_{\infty},
\end{equation}
and also some $q_0>0$ and $x_0>0$ such that
\begin{equation}\label{e:lyap1b}
-q_0 c(x)/x + \int_{(0,1)} (r^{-q_0}-1) \bar k (x,\dd r) \leq 0
\qquad\text{ for all }x\leq x_0.
\end{equation}
Then, \eqref{E:main}, and hence the strong Malthusian behavior of Theorem~\ref{t:main}, holds, and $\lambda > d$.
\end{prop}



\begin{proof}
In~\cite[Proposition~3.4$\MK$(ii)]{BW}, it was shown that, if the Markov process $X$ is recurrent, and the case when $c(x)=dx$ is a linear function excluded, then $\lambda > \inf_{x>0}c(x)/x$, so that~\eqref{E:main} holds. The further conditions ensure that $X$ is recurrent, as pointed out in~\cite[Section~3.6]{BeFK}.
\end{proof}

This proposition is effective in the case $c(x) \sim dx$ as $x\to 0$ and $x\to\infty$, provided $c$ is not linear. If $c$ is linear, $c(x) = dx$, then $\lambda=d$, $h(x)=x$, and one readily checks that the martingale $W$ is actually constant, so this case is less interesting.

We will now find a weaker version of the conditions above, which is effective when $d=0$. In this direction, if we can show that $\lambda > 0$, then the simple condition
\begin{equation}\label{e:c0}
\lim_{x\to 0+}\frac{c(x)}{x} = \lim_{x\to \infty}\frac{c(x)}{x} =0
\end{equation}
implies the strong Malthusian behavior of Theorem~\ref{t:main}.

Recall first that~\cite[Theorem~1]{DoumGab} already gives sufficient conditions for the strict positivity of the leading eigenvalue in the eigenproblem~\eqref{E:eigenp}. We can use our probabilistic methods to find alternative conditions for this.

For the sake of simplicity, we focus on the situation when fissions are binary (which puts $\Zz$ in the class of Markovian growth-fragmentations defined in~\cite{BeGF}). However, it is immediate to adapt the argument to the general case.

Assume that, for all $x>0$, $\kappa(x, \dd \pp)$ is supported by the set of binary mass partitions $\pp=(1-r,r, 0, \ldots)$ with $r\in(0,1/2]$. It is then more convenient to represent the fission kernel $\kappa$ by a probability kernel $\varrho(x,\dd r)$ on $(0,1/2]$, such that for all functionals $g\geq 0$ on $\Pp$,
\[
\int_{\Pp} g(\pp) \kappa(x, \dd \pp) = \int_{(0,1/2]} g(1-r,r, 0, \ldots) \varrho(x,\dd r).
\]
In particular, there is the identity
\[
\int_{(0,1)} f(r) \bar k (x,\dd r)= B(x)\int_{(0,1/2]} ((1-r) f(1-r)+ r f(r)) \varrho(x,\dd r).
\]

\begin{prop}\label{P1}
Assume conditions~\eqref{e:assump-c}--\eqref{e:assump-kappa}, \eqref{e:condc}--\eqref{e:assump} and, in the notation above, that there exist $q_{\infty}, x_{\infty}>0$ such that
\begin{equation}\label{e:Brho1}
q_{\infty} c(x)/x + B(x) \int_{(0,x1/2]} (r^{q_{\infty}}-1) \varrho (x, \dd r) \leq 0 \quad\text{ for all }x\geq x_{\infty},
\end{equation}
and $q_0, x_0>0$ such that
\begin{equation}\label{e:Brho2}
-q_0 c(x)/x + B(x) \int_{(0,1/2]} ((1-r)^{-q_0}-1) \varrho(x,\dd r) \leq 0 \quad\text{ for all }x\leq x_0.
\end{equation}
Then, the Malthus exponent $\lambda$ is positive. If moreover~\eqref{e:c0} holds, then condition~\eqref{E:main}, and hence the strong Malthusian behavior of Theorem~\ref{t:main}, holds.
\end{prop}

\begin{proof}
Let $a \in (x_0,x_\infty)$. By the definition of the Malthus exponent in Section~\ref{section4} and the right-continuity of the function $L_{a,a}$, we see that $\lambda>0$ if and only if $L_{a,a}(0)\in(1,\infty]$. We thus have to check that
\[
\EE_a\left[\Ee_{H(a)}, H(a)<\infty\right]>1,
\]
that is, thanks to the many-to-one formula of Lemma~\ref{L1'} (with $f=\mathbf{1}$), that
\[
\EE_a[\crochet{\Zz_{H(a)}, \mathbf{1}}]>1,
\]
where $\mathbf{1}$ is the constant function with value $1$. A fortiori, it suffices to check that $\crochet{\Zz_{H(a)}, \mathbf{1}}\geq 1$ $\PP_a$-a.s., and that this inequality is strict with positive $\PP_a$-probability. In words, we freeze individuals at their first return time to $a$; it is easy to construct an event with positive probability on which there are two or more frozen individuals, so we only need to verify that we get at least one frozen individual $\PP_a$-a.s. The latter point is not obvious, since it could in principle occur that no descendant of an individual started from $a$ succeeds in to returning to $a$.

In this direction, we focus on a specific ancestral trajectory, say $X^*$, which is defined as follows. Recall that any trajectory is driven by the flow velocity~\eqref{e:ode} between consecutive times of downward jumps, so we only need to explain how we select daughters at fission events. When a fission event occurs at a time $t$ with $X^*_{t-}=x^*>a$, producing two daughters, say $rx^*$ and $(1-r)x^*$ for some $r\in(0,1/2]$, then we choose the smallest daughter, i.e. $X^*_t=rx^*$, whereas if $x^*<a$ then we choose the largest daughter, i.e. $X^*_t=(1-r)x^*$. The process $X^*$ is then Markovian with infinitesimal generator
\[
{\mathcal G}^*f(x)=c(x) f'(x) +B(x) \int_{(0,1/2]} (\Indic{x>a} f(rx) + \Indic{x<a}f((1-r)x) -f(x))\varrho(x,\dd r).
\]

We now outline a proof that $X^*$ is point-recurrent. Since the process has only negative jumps, it is sufficient to show that $X^*_t \to \infty$ and $X^*_t \to 0$, as $t\to\infty$, are both impossible. For the former, consider starting the process at $x>x_\infty$ and killing it upon passage below $x_\infty$. Denote this process by $X^\circledast$ and its generator by
\[
{\mathcal G}^\circledast f(x)=c(x) f'(x) +B(x) \int_{(0,1/2]} (\Indic{rx > x_\infty} f(rx) -f(x))\varrho(x,\dd r).
\]
(The dependence on $a$ in the integral vanishes since $x_\infty>a$.) Now, let $V(x) = x^{q_\infty}$, for $x\ge x_\infty$. The conditions in the statement imply that $\mathcal{G}^\circledast V \le 0$, so $V(X^\circledast)$ is a supermartingale. This ensures that $X^\circledast$ cannot converge to $+\infty$, and indeed the same for $X^*$ itself. To show $X^*$ cannot converge to $0$, we start it at $x<x_0$ and kill it upon passing above $x_0$, and follow the same argument with $V(x) = x^{-q_0}$.

To conclude, we have shown that $X^*$ is point-recurrent, and therefore $\PP_a$-almost surely hits $a$. This shows that $\PP_a[\crochet{\Zz_{H(a)}, \mathbf{1}}\geq 1]=1$, and completes the proof of Proposition~\ref{P1}.
\end{proof}
\goodbreak


The conditions in Propositions~\ref{P0} and~\ref{P1} express that the fragmentation is stronger than growth near $\infty$, and the reverse near $0$. Proposition~\ref{P0} describes this comparison in terms of the Markov process $X$, which can be selected from the process $\mathbb{Z}$ by making a size-biased pick from the offspring at each branching event; that is, from offspring of sizes $rx$ and $(1-r)x$, following the former with probability $r$ and the latter with probability $1-r$. On the other hand, Proposition~\ref{P1} describes this comparison for the process $X^*$, where we pick from the offspring more carefully in order to follow a line of descent which is more likely to stay close to the point $a$. This accounts for the improvement in conditions between~\cite{BeFK} and this work.

Finally, we conclude with a discussion of the self-similar case, where our conditions can be simplified further. We say that the fragmentation is \emph{self-similar} if the relative sizes of offspring do not depend on the size of the parent; that is, $\kappa(x,\dd \pp) = \kappa(\dd \pp)$. We define a measure $\sigma$ on $(0,1)$ by
\[
\int_{(0,1)} f(r) \, \sigma(\dd r) = \int_{\Pp} \sum_{i=1}^\infty f(p_i) \, \kappa(\dd \pp),
\]
and note that then the measure $\bar{k}$, defined in terms of $\kappa$ on page~\pageref{s:kbar}, satisfies
\[
\bar{k}(x,\dd r) = B(x) \bar{\sigma}(\dd r),
\]
where $\bar{\sigma}(\dd r) = r\sigma(\dd r)$ is a probability measure. If this is the case, then conditions~\eqref{e:lyap1a} and~\eqref{e:lyap1b} of Proposition~\ref{P0} can be rewritten as
\begin{align}
\frac{x B(x)}{c(x)} &\ge \frac{q_\infty}{\int_{(0,1)} (1-r^{q_\infty})\, \bar{\sigma}(\dd r)},
\quad \;x \ge x_\infty,\label{e:Bc10} \\
\intertext{and}
\frac{x B(x)}{c(x)} &\le \frac{q_0}{\int_{(0,1)} (r^{-q_0}-1)\, \bar{\sigma}(\dd r)},
\quad x \le x_0.\label{e:Bc20}
\end{align}
The comparable conditions in~\cite{DoumGab} would be (13) and (12), respectively, which state that $\lim_{x\to\infty}x B(x)/c(x) = \infty$ and $B/c \in L^1((0,b), \dd x)$ for some $b>0$. Our conditions are rather weaker than these, though of course in other respects~\cite{DoumGab} is more general.

Turning to Proposition~\ref{P1}, self-similarity means that we have $\varrho(x,\dd r) = \varrho(\dd r)$, and $\sigma$ can be expressed as $\sigma(\dd r) = \varrho(\dd r) + \varrho(1-\dd r)$. The conditions~\eqref{e:Brho1} and~\eqref{e:Brho2} of Proposition~\ref{P1} become
\begin{align}
\frac{x B(x)}{c(x)} &\ge \frac{q_\infty}{\int_{(0,1/2]} (1-r^{q_\infty})\, \varrho(\dd r)},  \qquad\quad\;\;\; x\ge x_\infty,\label{e:Bc1}\\
\intertext{and}
\frac{x B(x)}{c(x)} &\le \frac{q_0}{\int_{(0,1/2]} ((1-r)^{-q_0}-1)\, \varrho(\dd r)}, \quad x\le x_0.\label{e:Bc2}
\end{align}

To demonstrate these conditions, consider the simple case of binary fission with uniform size repartition. In this case, $\bar \sigma(\dd r) = 2r \dd r $, and the right-hand sides of~\eqref{e:Bc10} and~\eqref{e:Bc20} become $q_{\infty}+2$ and $2-q_0$, respectively. Hence, the conditions~\eqref{e:lyap1a} and~\eqref{e:lyap1b} of Proposition~\ref{P0} hold if there exist $0<x_0<x_{\infty}$ and $\beta_0<2<\beta_{\infty}$ such that
\[
\frac{x B(x)}{c(x)} \geq \beta_{\infty} \quad \text{for all } x \ge x_\infty
\quad \text{and} \quad
\frac{x B(x)}{c(x)} \leq \beta_0 \quad \text{for all }x \le x_0.
\]
In turn, Proposition~\ref{P1} improves upon the preceding bounds. Specifically, one has $\varrho(\dd r) = 2\cdot {\mathbf 1}_{(0,1/2]}(r) \dd r$, and the right-hand side of~\eqref{e:Bc1} becomes $1-\frac{2^{-q_\infty}}{q_\infty+1}$, which is an increasing function whose limit as $q_\infty\to 0$ is $\frac{1}{1+\log 2} \approx 0.59$. The right-hand side of~\eqref{e:Bc2} becomes $\frac{2}{1-q_0}(1-2^{q_0-1})-1$ (for $q_0\ne 1$), a decreasing function with limit $\frac{1}{1-\log 2} \approx 3.3$ as $q_0\to 0$. Hence, the conditions of Proposition~\ref{P1} hold if there exist $0<x_0<x_{\infty}$, $\beta'_{\infty}>\frac{1}{1+\log 2}$ and $\beta'_0<\frac{1}{1-\log 2}$, such that
\[
\frac{x B(x)}{c(x)} \geq \beta'_{\infty} \quad \text{for all } x \ge x_\infty
\quad \text{and} \quad
\frac{x B(x)}{c(x)} \leq \beta'_0 \quad \text{for all }x \le x_0.
\]
Again, these are weaker than~\cite[Conditions (12) and (13)]{DoumGab}.

\subsection*{Acknowledgments}
We thank two anonymous referees for their careful and thorough reviews.


\begin{thebibliography}{MMMMMM}
\bibitem[AH76]{AsHer} Asmussen, Soren; Hering, Heinrich Strong limit theorems for general supercritical branching processes with applications to branching diffusions, Z. Wahrscheinlichkeitstheor. Verw. Geb., Volume 36 (1976) no. 3, pp. 195-212 \href{https://ahl.centre-mersenne.org/item/AHL_2020__3__795_0/#AsHer}{Publisher reference}.
\bibitem[BBCK18]{BBCK} Bertoin, Jean; Budd, Timothy; Curien, Nicolas; Kortchemski, Igor Martingales in self-similar growth-fragmentations and their connections with random planar maps, Probab. Theory Relat. Fields, Volume 172 (2018) no. 3-4, pp. 663-724 \href{https://ahl.centre-mersenne.org/item/AHL_2020__3__795_0/#BBCK}{Publisher reference}.
\bibitem[BBH + 15]{BBHHR} Berestycki, Julien; Brunet, Éric; Harris, John W.; Harris, Simon C.; Roberts, Matthew I. Growth rates of the population in a branching Brownian motion with an inhomogeneous breeding potential, Stochastic Processes Appl., Volume 125 (2015) no. 5, pp. 2096-2145 \href{https://ahl.centre-mersenne.org/item/AHL_2020__3__795_0/#BBHHR}{Publisher reference}.
\bibitem[BCG13]{BCG-fine} Balagué, Daniel; Cañizo, José A.; Gabriel, Pierre Fine asymptotics of profiles and relaxation to equilibrium for growth-fragmentation equations with variable drift rates, Kinet. Relat. Models, Volume 6 (2013) no. 2, pp. 219-243 \href{https://ahl.centre-mersenne.org/item/AHL_2020__3__795_0/#BCG-fine}{Publisher reference}.
\bibitem[BCG20]{BanCloGab} Bansaye, Vincent; Cloez, Bertrand; Gabriel, Pierre Ergodic Behavior of Non-conservative Semigroups via Generalized Doeblin’s Conditions, Acta Appl. Math., Volume 166 (2020), pp. 29-72 \href{https://ahl.centre-mersenne.org/item/AHL_2020__3__795_0/#BanCloGab}{Publisher reference}.
\bibitem[BCGM19]{BCGM-qs} Bansaye, Vincent; Cloez, Bertrand; Gabriel, Pierre; Marguet, Aline A non-conservative Harris’ ergodic theorem (2019) (\url{https://arxiv.org/abs/1903.03946v1}) \href{https://ahl.centre-mersenne.org/item/AHL_2020__3__795_0/#BCGM-qs}{Publisher reference}.
\bibitem[BDMT11]{BDMT} Bansaye, Vincent; Delmas, Jean-François; Marsalle, Laurence; Tran, Viet Chi Limit theorems for Markov processes indexed by continuous time Galton–Watson trees, Ann. Appl. Probab., Volume 21 (2011) no. 6, pp. 2263-2314 \href{https://ahl.centre-mersenne.org/item/AHL_2020__3__795_0/#BDMT}{Publisher reference}.
\bibitem[Ber03]{Beresty} Berestycki, Julien Multifractal spectra of fragmentation processes, J. Stat. Phys., Volume 113 (2003) no. 3-4, pp. 411-430 \href{https://ahl.centre-mersenne.org/item/AHL_2020__3__795_0/#Beresty}{Publisher reference}.
\bibitem[Ber06]{Ber-frag} Bertoin, Jean Random fragmentation and coagulation processes, Cambridge Studies in Advanced Mathematics, 102, Cambridge University Press, 2006, viii+280 pages \href{https://ahl.centre-mersenne.org/item/AHL_2020__3__795_0/#Ber-frag}{Publisher reference}.
\bibitem[Ber17]{BeGF} Bertoin, Jean Markovian growth-fragmentation processes, Bernoulli, Volume 23 (2017) no. 2, pp. 1082-1101 \href{https://ahl.centre-mersenne.org/item/AHL_2020__3__795_0/#BeGF}{Publisher reference}.
\bibitem[Ber19]{BeFK} Bertoin, Jean On a Feynman–Kac approach to growth-fragmentation semigroups and their asymptotic behaviors, J. Funct. Anal., Volume 277 (2019) no. 11, 108270, p. 29 \href{https://ahl.centre-mersenne.org/item/AHL_2020__3__795_0/#BeFK}{Publisher reference}.
\bibitem[BG20]{BerGab2} Bernard, Étienne; Gabriel, Pierre Asynchronous exponential growth of the growth-fragmentation equation with unbounded fragmentation rate, J. Evol. Equ., Volume 20 (2020), pp. 375-401 \href{https://ahl.centre-mersenne.org/item/AHL_2020__3__795_0/#BerGab2}{Publisher reference}.
\bibitem[Big92]{Biggins92} Biggins, J. D. Uniform convergence of martingales in the branching random walk, Ann. Probab., Volume 20 (1992) no. 1, pp. 137-151 \href{https://ahl.centre-mersenne.org/item/AHL_2020__3__795_0/#Biggins92}{Publisher reference}.
\bibitem[BK04]{BigKyp} Biggins, J. D.; Kyprianou, Andreas E. Measure change in multitype branching, Adv. Appl. Probab., Volume 36 (2004) no. 2, pp. 544-581 \href{https://ahl.centre-mersenne.org/item/AHL_2020__3__795_0/#BigKyp}{Publisher reference}.
\bibitem[BR05]{BeRou} Bertoin, Jean; Rouault, Alain Discretization methods for homogeneous fragmentations, J. London Math. Soc. (2), Volume 72 (2005) no. 1, pp. 91-109 \href{https://ahl.centre-mersenne.org/item/AHL_2020__3__795_0/#BeRou}{Publisher reference}.
\bibitem[BW18]{BW} Bertoin, Jean; Watson, Alexander R. A probabilistic approach to spectral analysis of growth-fragmentation equations, J. Funct. Anal., Volume 274 (2018) no. 8, pp. 2163-2204 \href{https://ahl.centre-mersenne.org/item/AHL_2020__3__795_0/#BW}{Publisher reference}.
\bibitem[Cav20]{Cavalli} Cavalli, Benedetta On a family of critical growth-fragmentation semigroups and refracted Lévy processes, Acta Appl. Math., Volume 166 (2020) no. 1, pp. 161-186 \href{https://ahl.centre-mersenne.org/item/AHL_2020__3__795_0/#Cavalli}{Publisher reference}.
\bibitem[CDP18]{CalvoDoumicPerthame} Calvo, Juan G.; Doumic, Marie; Perthame, Benoît Long-Time Asymptotics for Polymerization Models, Commun. Math. Phys., Volume 363 (2018) no. 1, pp. 111-137 \href{https://ahl.centre-mersenne.org/item/AHL_2020__3__795_0/#CalvoDoumicPerthame}{Publisher reference}.
\bibitem[Cha91]{Cha-stop} Chauvin, Brigitte Product martingales and stopping lines for branching Brownian motion, Ann. Probab., Volume 19 (1991) no. 3, pp. 1195-1205 \href{https://ahl.centre-mersenne.org/item/AHL_2020__3__795_0/#Cha-stop}{Publisher reference}.
\bibitem[CHHK19]{CHHK} Cox, Alexander M. G.; Harris, Simon C.; Horton, Emma L.; Kyprianou, Andreas E. Multi-species Neutron Transport Equation, J. Stat. Phys., Volume 176 (2019) no. 2, pp. 425-455 \href{https://ahl.centre-mersenne.org/item/AHL_2020__3__795_0/#CHHK}{Publisher reference}.
\bibitem[Clo17]{Cloez} Cloez, Bertrand Limit theorems for some branching measure-valued processes, Adv. Appl. Probab., Volume 49 (2017) no. 2, pp. 549-580 \href{https://ahl.centre-mersenne.org/item/AHL_2020__3__795_0/#Cloez}{Publisher reference}.
\bibitem[CRY17]{CRY} Chen, Zhen-Qing; Ren, Yan-Xia; Yang, Ting Law of large numbers for branching symmetric Hunt processes with measure-valued branching rates, J. Theor. Probab., Volume 30 (2017) no. 3, pp. 898-931 \href{https://ahl.centre-mersenne.org/item/AHL_2020__3__795_0/#CRY}{Publisher reference}.
\bibitem[CS07]{ChYu} Chen, Zhen-Qing; Shiozawa, Yuichi Limit theorems for branching Markov processes, J. Funct. Anal., Volume 250 (2007) no. 2, pp. 374-399 \href{https://ahl.centre-mersenne.org/item/AHL_2020__3__795_0/#ChYu}{Publisher reference}.
\bibitem[CV16]{ChVil} Champagnat, Nicolas; Villemonais, Denis Exponential convergence to quasi-stationary distribution and Q-process, Probab. Theory Relat. Fields, Volume 164 (2016) no. 1-2, pp. 243-283 \href{https://ahl.centre-mersenne.org/item/AHL_2020__3__795_0/#ChVil}{Publisher reference}.
\bibitem[CV20]{CV-unbounded} Champagnat, Nicolas; Villemonais, Denis Practical criteria for R-positive recurrence of unbounded semigroups, Electron. Commun. Probab., Volume 25 (2020), 6 \href{https://ahl.centre-mersenne.org/item/AHL_2020__3__795_0/#CV-unbounded}{Publisher reference}.
\bibitem[Dad17]{Dadoun} Dadoun, Benjamin Asymptotics of self-similar growth-fragmentation processes, Electron. J. Probab., Volume 22 (2017), 27, p. 30 \href{https://ahl.centre-mersenne.org/item/AHL_2020__3__795_0/#Dadoun}{Publisher reference}.
\bibitem[DB92]{DB-bpre} D’Souza, J. C.; Biggins, J. D. The supercritical Galton–Watson process in varying environments, Stoch. Proc. Appl., Volume 42 (1992) no. 1, pp. 39-47 \href{https://ahl.centre-mersenne.org/item/AHL_2020__3__795_0/#DB-bpre}{Publisher reference}.
\bibitem[DDGW18]{DDGW} D{\fontspec{DejaVu Serif}ȩ}biec, Tomasz; Doumic, Marie; Gwiazda, Piotr; Wiedemann, Emil Relative entropy method for measure solutions of the growth-fragmentation equation, SIAM J. Math. Anal., Volume 50 (2018) no. 6, pp. 5811-5824 \href{https://ahl.centre-mersenne.org/item/AHL_2020__3__795_0/#DDGW}{Publisher reference}.
\bibitem[DJG10]{DoumGab} Doumic Jauffret, Marie; Gabriel, Pierre Eigenelements of a general aggregation-fragmentation model, Math. Models Methods Appl. Sci., Volume 20 (2010) no. 5, pp. 757-783 \href{https://ahl.centre-mersenne.org/item/AHL_2020__3__795_0/#DoumGab}{Publisher reference}.
\bibitem[DM04]{Delmoral} Del Moral, Pierre Feynman–Kac formulae. Genealogical and interacting particle systems with applications, Probability and its Applications, Springer, 2004, xviii+555 pages \href{https://ahl.centre-mersenne.org/item/AHL_2020__3__795_0/#Delmoral}{Publisher reference}.
\bibitem[EHK10]{EnHaKy} Engländer, János; Harris, Simon C.; Kyprianou, Andreas E. Strong law of large numbers for branching diffusions, Ann. Inst. Henri Poincaré, Probab. Stat., Volume 46 (2010) no. 1, pp. 279-298 \href{https://ahl.centre-mersenne.org/item/AHL_2020__3__795_0/#EnHaKy}{Publisher reference}.
\bibitem[EN00]{EN00} Engel, Klaus-Jochen; Nagel, Rainer One-parameter semigroups for linear evolution equations, Graduate Texts in Mathematics, 194, Springer, 2000, xxii+586 pages \href{https://ahl.centre-mersenne.org/item/AHL_2020__3__795_0/#EN00}{Publisher reference}.
\bibitem[Esc20]{Esco} Escobedo, Miguel On the non existence of non negative solutions to a critical Growth-Fragmentation Equation, Ann. Fac. Sci. Toulouse, Math., Volume 29 (2020) no. 1, pp. 177-220 \href{https://ahl.centre-mersenne.org/item/AHL_2020__3__795_0/#Esco}{Publisher reference}.
\bibitem[GHH07]{GiHaHa} Git, Y.; Harris, John W.; Harris, Simon C. Exponential growth rates in a typed branching diffusion, Ann. Appl. Probab., Volume 17 (2007) no. 2, pp. 609-653 \href{https://ahl.centre-mersenne.org/item/AHL_2020__3__795_0/#GiHaHa}{Publisher reference}.
\bibitem[HH09]{HaHa} Hardy, Robert; Harris, Simon C. A spine approach to branching diffusions with applications to L p -convergence of martingales, Séminaire de Probabilités XLII (Catherine, Donati-Martin; al, eds.) (Lecture Notes in Mathematics), Volume 1979, Springer, 2009, pp. 281-330 \href{https://ahl.centre-mersenne.org/item/AHL_2020__3__795_0/#HaHa}{Publisher reference}.
\bibitem[HHK16]{HaHeKy} Harris, Simon C.; Hesse, Marion; Kyprianou, Andreas E. Branching Brownian motion in a strip: survival near criticality, Ann. Probab., Volume 44 (2016) no. 1, pp. 235-275 \href{https://ahl.centre-mersenne.org/item/AHL_2020__3__795_0/#HaHeKy}{Publisher reference}.
\bibitem[HHK19]{HHK} Harris, Simon C.; Horton, Emma L.; Kyprianou, Andreas L. Stochastic Methods for the Neutron Transport Equation II: Almost sure growth (2019) (\url{https://arxiv.org/abs/1901.00220,} to appear in The Annals of Applied Probability) \href{https://ahl.centre-mersenne.org/item/AHL_2020__3__795_0/#HHK}{Publisher reference}.
\bibitem[HKV18]{HKV} Horton, Emma; Kyprianou, Andreas; Villemonais, Denis Stochastic Methods for the Neutron Transport Equation I: Linear Semigroup asymptotics (2018) (\url{https://arxiv.org/abs/1810.01779,} to appear in The Annals of Applied Probability) \href{https://ahl.centre-mersenne.org/item/AHL_2020__3__795_0/#HKV}{Publisher reference}.
\bibitem[Jag89]{Jagers} Jagers, Peter General branching processes as Markov fields, Stoch. Proc. Appl., Volume 32 (1989) no. 2, pp. 183-212 \href{https://ahl.centre-mersenne.org/item/AHL_2020__3__795_0/#Jagers}{Publisher reference}.
\bibitem[JN84]{JN84} Jagers, Peter; Nerman, Olle The growth and composition of branching populations, Adv. Appl. Probab., Volume 16 (1984) no. 2, pp. 221-259 \href{https://ahl.centre-mersenne.org/item/AHL_2020__3__795_0/#JN84}{Publisher reference}.
\bibitem[KP76]{KaP} Kahane, Jean-Pierre; Peyrière, Jacques Sur certaines martingales de Benoit Mandelbrot, Adv. Math., Volume 22 (1976) no. 2, pp. 131-145 \href{https://ahl.centre-mersenne.org/item/AHL_2020__3__795_0/#KaP}{Publisher reference}.
\bibitem[Mar19]{Mar-lln} Marguet, Aline A law of large numbers for branching Markov processes by the ergodicity of ancestral lineages, ESAIM Probab. Stat., Volume 23 (2019), pp. 638-661 \href{https://ahl.centre-mersenne.org/item/AHL_2020__3__795_0/#Mar-lln}{Publisher reference}.
\bibitem[MS16]{MS16} Mischler, Stéphane; Scher, Joshua Spectral analysis of semigroups and growth-fragmentation equations, Ann. Inst. H. Poincaré Anal. Non Linéaire, Volume 33 (2016) no. 3, pp. 849-898 \href{https://ahl.centre-mersenne.org/item/AHL_2020__3__795_0/#MS16}{Publisher reference}.
\bibitem[Ner81]{Nerman} Nerman, Olle On the convergence of supercritical general (C-M-J) branching processes, Z. Wahrscheinlichkeitstheor. Verw. Geb., Volume 57 (1981) no. 3, pp. 365-395 \href{https://ahl.centre-mersenne.org/item/AHL_2020__3__795_0/#Nerman}{Publisher reference}.
\bibitem[Per07]{Per07} Perthame, Benoît Transport equations in biology, Frontiers in Mathematics, Birkhäuser, 2007 \href{https://ahl.centre-mersenne.org/item/AHL_2020__3__795_0/#Per07}{Publisher reference}.
\bibitem[Pey74]{Pey} Peyrière, Jacques Turbulence et dimension de Hausdorff, C. R. Acad. Sci., Paris, Sér. A, Volume 278 (1974), pp. 567-569 \href{https://ahl.centre-mersenne.org/item/AHL_2020__3__795_0/#Pey}{Publisher reference}.
\bibitem[Sav69]{Savits} Savits, Thomas H. The explosion problem for branching Markov process, Osaka J. Math., Volume 6 (1969) no. 2, pp. 375-395 \href{https://ahl.centre-mersenne.org/item/AHL_2020__3__795_0/#Savits}{Publisher reference}.
\bibitem[Shi08]{Shio} Shiozawa, Yuichi Exponential growth of the numbers of particles for branching symmetric \(\alpha\)-stable processes, J. Math. Soc. Japan, Volume 60 (2008) no. 1, pp. 75-116 \href{https://ahl.centre-mersenne.org/item/AHL_2020__3__795_0/#Shio}{Publisher reference}.
\bibitem[Shi15]{Shi} Shi, Zhan Branching random walks. École d’Été de Probabilités de Saint-Flour XLII – 2012, Lecture Notes in Mathematics, 2151, Springer, 2015, x+133 pages \href{https://ahl.centre-mersenne.org/item/AHL_2020__3__795_0/#Shi}{Publisher reference}.
\bibitem[Shi20]{ShiQ-ou} Shi, Quan A growth-fragmentation model related to Ornstein–Uhlenbeck type processes, Ann. Inst. H. Poincaré Probab. Statist., Volume 56 (2020) no. 1, pp. 580-611 \href{https://ahl.centre-mersenne.org/item/AHL_2020__3__795_0/#ShiQ-ou}{Publisher reference}.
\end{thebibliography}
\end{document}
